\begin{textsample}{POS Dim 4 – 2021\_c6 – Score 9.00 – 2021\_c6/enade\_2021\_geo\_b7744036e1.txt}  \label{ex:f4_pos_061}
CAPÍTULO 6 ANÁLISE TÉCNICA DA PROVA Este capítulo tem por objetivo apresentar o desempenho dos estudantes concluintes de Geografia ( Licenciatura/Bacharelado ) no Enade/2021.
Para isso, foram calculadas as estatísticas básicas da prova em seu todo ( seção 6.1.1 ), bem como as estatísticas relacionadas aos Componentes de Formação Geral ( seção 6.1.2 ) e de Conhecimento Específico da Área ( seção 6.1.3 ).
Dadas as suas características, foram analisadas, em separado, as questões objetivas ( seção 6.2 ) e as questões discursivas ( seção 6.3 ).
Tomando -se como base as duas Questões Discursivas do Componente de Formação Geral, nas seções 6.3.1.6 a 6.3.1.8, são apresentados os resultados e os comentários sobre a correção das respostas em relação à Língua Portuguesa.
Para as questões objetivas, foram disponibilizados os Índices de Facilidade e de Discriminação Ponto-Bisserial, também em separado, para os Componentes de Formação Geral ( seção 6.2.1 e 6.2.2 ) e de Conhecimento Específico ( seção 6.2.3 e 6.2.4 ).
De cada componente, uma das questões foi escolhida para exemplificar a análise gráfica, relacionando -se as alternativas escolhidas pelos estudantes ( inclusive o gabarito ) com o número de acertos no componente.
Nos Anexos IA e IB, apresenta -se a íntegra da análise gráfica para todas as questões objetivas, respectivamente à Licenciatura e ao Bacharelado.
Para cada uma das questões discursivas, os conteúdos dos tipos mais comuns de respostas dos estudantes são apresentados e comparados com o padrão de resposta esperado ( ver Anexos IXA e IXB com o padrão de respostas ).
Nas tabelas, constam as seguintes estatísticas das notas19 : média do desempenho na prova, erro padrão da média, desvio padrão, nota mínima, mediana e nota máxima para cada um de seus componentes.
Tais estatísticas contemplam o total de estudantes concluintes da Área de Geografia ( Licenciatura/Bacharelado ) inscritos e presentes à prova do Enade/2021, tendo em vista agregações, ou por Grandes Regiões e o país como um todo, segundo Habilitação e Modalidade de Ensino, ou por Categoria Administrativa, Organização Acadêmica da IES.
Em relação aos gráficos de distribuição de notas, o intervalo considerado foi de 10 unidades, aberto à esquerda e fechado à direita, com exceção do primeiro intervalo, [ 0 ; 10 ], fechado em ambos os extremos.
Para os gráficos de distribuição das notas das questões discursivas, foram consideradas mais duas categorias : questão em branco20 e nota zero. 19 Uma definição dessas estatísticas pode ser encontrada no Glossário. 20 Nesse grupo estão incluídas também as respostas classificadas como nulas ou desconsideradas. 6.1 ESTATÍSTICAS BÁSICAS DA PROVA Esta seção apresenta estatísticas selecionadas e histogramas da nota geral ( 6.1.1 ) e de cada componente : Formação Geral ( 6.1.2 ) e Conhecimento Específico ( 6.1.3 ).
São também apresentadas estatísticas selecionadas de subpopulações, caracterizadas por Grande Região, Categoria Administrativa, Organização Acadêmica e Modalidade de Ensino. 6.1.1 ESTATÍSTICAS BÁSICAS GERAIS Na Tabela 6.1, são apresentadas as Estatísticas Básicas da prova, por Grande Região segundo Habilitação e Modalidade de Ensino, dos estudantes concluintes de Geografia ( Licenciatura/Bacharelado ).
As análises serão feitas primeiramente para os cursos de Licenciatura e depois para os de Bacharelado.
Cada Habilitação é desagregada internamente em cursos Presenciais e cursos a Distância.
A Média das notas da prova, como um todo, dos alunos de Licenciatura ( nas seções seguintes serão analisados os Componentes de Formação Geral e de Conhecimento Específico ), foi 40,8, sendo que os estudantes da região Norte obtiveram a Média mais baixa ( 36,0 ), e os da região Sudeste obtiveram a Média mais alta ( 44,8 ).
As demais Médias foram : 39,2, na região Nordeste, 39,8, na região Sul e 39,1, na região Centro-Oeste.
O Desvio padrão para o Brasil, como um todo, foi 14,3, sendo o maior Desvio padrão encontrado na região Sudeste ( 15,2 ), e o menor, na região Norte ( 12,3 ), indicando uma dispersão um pouco menor das notas dessa última região.
A região que obteve a maior nota Máxima foi a Sudeste ( 93,7 ), ao passo que a região que atingiu a menor nota Máxima foi a Norte ( 71,6 ).
A Mediana do Brasil, como um todo, foi 40,5, sendo a maior Mediana obtida na região Sudeste ( 45,2 ), e a menor obtida na região Norte ( 34,8 ).
A nota Mínima do Brasil foi obtida nas regiões Nordeste, Sudeste e Sul, onde foi zero.
As maiores notas Mínimas, entre as cinco Grandes Regiões, foram 5,4, na região Norte, e 1,9, na região Centro-Oeste.
Considerando -se as notas, como um todo dos alunos de Licenciatura segundo Grande Região, observa -se que não existe diferença estatisticamente significativa ao nível de 95 \% apenas entre as Médias obtidas nas regiões Nordeste, Sul e Centro-Oeste.
A Média das notas para os alunos de cursos Presenciais de Licenciatura, foi 42,5, sendo que os alunos da região Norte obtiveram a Média mais baixa ( 36,0 ), e os da região Sudeste obtiveram a Média mais alta ( 47,9 ).
As demais Médias foram : 40,2 na região Nordeste, 47,2 na região Sul e 39,1 na região Centro-Oeste.
O Desvio padrão das Licenciaturas Presenciais foi 14,3, sendo o menor obtido na região Norte 12,3 e o maior na região Sudeste ( 14,7 ).
A maior Mediana foi a da região Sudeste ( 49,3 ) e a menor foi a da região Norte ( 34,8 ).
A maior nota Máxima foi 86,6, obtida na região Sul, registrada como a nota Máxima de todas as Licenciaturas.
As notas Mínimas foram diferentes de zero nas regiões Norte ( 5,4 ), Sul ( 7,0 ) e Centro-Oeste ( 6,4 ). região Sudeste ( 14,7 ).
A maior Mediana foi a da região Sudeste ( 49,3 ) e a menor foi a da região Norte ( 34,8 ).
A maior nota Máxima foi 86,6, obtida na região Sul, registrada como a nota Máxima de todas as Licenciaturas.
As notas Mínimas foram diferentes de zero nas regiões Norte ( 5,4 ), Sul ( 7,0 ) e Centro-Oeste ( 6,4 ).
Na Modalidade de Educação Presencial de Licenciatura as diferenças entre Médias obtidas nas regiões Nordeste e Centro-Oeste não são estatisticamente significativas, assim como entre as regiões Sudeste e Sul.
A região Norte foi a única que não apresentou cursos de Licenciatura avaliados na modalidade Educação a Distância.
A Média das notas da prova dos alunos de cursos a Distância de Licenciatura, foi 38,6, mais baixa do que os alunos de cursos Presenciais ( 42,5 ) sendo que os alunos da região Nordeste obtiveram a Média mais baixa ( 33,4 ), e os da região Sudeste obtiveram a Média mais alta ( 41,1 ).
As demais Médias foram : 38,2 na região Sul e 39,4 na região Centro-Oeste.
O Desvio padrão das Licenciaturas a Distância foi 14,0, sendo o menor obtido na região Nordeste ( 12,4 ) e o maior na região Sudeste ( 14,9 ).
A maior Mediana foi a da região Sudeste ( 40,3 ) e a menor foi a da região Nordeste ( 33,1 ).
A maior nota Máxima foi 93,7, obtida na região Sudeste, registrada como a nota Máxima de todas as Licenciaturas.
As notas Mínimas foram : 6,5 na região Nordeste, 1,9 nas regiões Sudeste e Centro-Oeste, e zero na região Sul.
Na Modalidade de Educação a Distância de Licenciatura as diferenças entre Médias obtidas na região Nordeste e nas regiões Sudeste, Sul e Centro-Oeste são estatisticamente significativas, assim como entre as regiões Sudeste e Sul.
Cumpre notar que ainda que as notas dos cursos Presenciais sejam comparáveis com as dos cursos a Distância dentro de uma mesma Habilitação, não é possível comparar as notas obtidas pelos alunos de Licenciatura e de Bacharelado, já que as provas são, em parte, diferentes.
A Média das notas da prova, como um todo dos alunos de Bacharelado ( nas seções seguintes serão analisados os Componentes de Formação Geral e de Conhecimento Específico ), foi 43,2, sendo que os estudantes da região Norte obtiveram a Média mais baixa ( 37,7 ), e os da região Sul obtiveram a Média mais alta ( 45,1 ).
As demais Médias foram : 41,4, na região Nordeste, 44,9, na região Sudeste e 38,7 na região Centro-Oeste.
O Desvio padrão para o Brasil, como um todo, foi 13,2, sendo o maior Desvio padrão encontrado na região Norte ( 13,7 ), e o menor, na região Nordeste ( 12,4 ), indicando uma dispersão um pouco menor das notas dessa última região.
A região que obteve a maior nota Máxima foi a Nordeste ( 76,3 ), ao passo que a região que atingiu a menor nota Máxima foi a Centro-Oeste ( 68,4 ).
A Mediana do Brasil, como um todo, foi 43,4, sendo a maior Mediana obtida na região Sudeste ( 45,2 ), e a menor obtida na região Norte ( 37,3 ).
A nota Mínima do Brasil foi obtida nas regiões Sudeste e Centro-Oeste, onde foi zero.
A maior nota Mínima, entre as cinco Grandes Regiões, foi 8,3, na região Norte ; nas demais regiões as notas Mínimas foram : 7,7 na região Nordeste e 4,5 na região Sul.
Considerando -se as notas de todo o Brasil dos alunos de Bacharelado, segundo Grande Região, observa -se que não existe diferença estatisticamente significativa ao nível de 95 \% entre as Médias, obtidas nas regiões Sudeste e Sul, assim como entre as regiões Norte, Nordeste e Centro-Oeste.
A Média das notas da prova como um todo para os alunos de cursos Presenciais de Bacharelado, foi 43,5, sendo que os alunos da região Norte obtiveram a Média mais baixa ( 37,7 ), e os da região Sul obtiveram a Média mais alta ( 45,9 ).
As demais Médias foram : 41,4 na região Nordeste, 45,3 na região Sudeste e 38,7 na região Centro-Oeste.
A maior nota Máxima dos cursos Presenciais de Bacharelado foi a da região Nordeste ( 76,3 ) e a menor foi a da região Centro-Oeste ( 68,4 ).
A maior Mediana foi obtida na região Sul ( 46,0 ), e a menor, na Norte ( 37,3 ).
As notas Mínimas nos cursos Presenciais foram nas regiões Norte ( 8,3 ), Nordeste ( 7,7 ), Sudeste ( 0,0 ), Sul ( 4,5 ) e Centro-Oeste ( 0,0 ).
Considerando -se as notas segundo Grande Região, não existe diferença estatisticamente significativa ao nível de 95 \% entre as Médias, obtidas nas regiões Sudeste e Sul, assim como entre as regiões Norte, Nordeste e Centro-Oeste.
Somente nas regiões Sudeste e Sul foram oferecidos cursos a Distância de Bacharelado.
Respectivamente, para o Brasil, a região Sudeste e a região Sul as Médias foram 38,9, 36,0 e 40,4 ; o Desvio padrão foi 13,6, 14,6 e 12,8 ; a nota Mínima foi 9,2, 9,2 e 14,7 ; a Mediana foi 38,0, 36,4 e 39,8 ; e a nota Máxima foi 74,6, 74,6 e 69,2.
Não existe diferença estatisticamente significativa ao nível de 95 \% entre as Médias, obtidas nas regiões Sudeste e Sul.
O comportamento das notas dos estudantes de todo o Brasil pode ser observado no Gráfico 6.1, que apresenta um histograma com a distribuição das mesmas, segundo a Habilitação.
Essa é uma distribuição unimodal para os alunos de cursos de Licenciatura e de Bacharelado, com modas no intervalo ( 30 ; 40 ] e ( 40 ; 50 ], respectivamente.
Na Tabela 6.2, são apresentadas informações referentes à Média da nota final, desagregadas por Categoria Administrativa, por Organização Acadêmica, segundo a Habilitação.
A Média das notas da prova, como um todo, dos alunos de Licenciatura, foi 40,8 ( Tabela 6.1 ).
Em relação a Categoria Administrativa, os estudantes de Licenciatura das IES Públicas obtiveram Média mais alta ( 41,7 ) que a Média nacional, e os das IES Privadas obtiveram Média mais baixa ( 39,1 ).
Observa -se que existe diferença estatisticamente significativa entre as Médias das notas das IES Públicas e as das Privadas.
A diferença entre as Médias das regiões Sudeste e Norte ( 8,8 ), a maior e a menor Média, é superior à diferença entre a das IES Públicas e Privadas ( 2,6 ), chegando a caracterizar uma maior diversidade regional do que administrativa.
O Desvio padrão para as IES Públicas ( 14,4 ) foi superior ao do Brasil, como um todo ( 14,3 ), indicando uma dispersão um pouco maior das notas nesta Categoria Administrativa.
No tocante à Organização Acadêmica, as Universidades ( 41,2 ) e os CEFET/IF ( 42,3 ) obtiveram Média mais alta que a nacional ( 40,8 ).
A Média dos Centros Universitários ( 39,1 ) e das Faculdades ( 34,7 ) foi menor do que a nacional.
Constata -se que não existe diferença 6.1.2 ESTATÍSTICAS BÁSICAS NO COMPONENTE DE FORMAÇÃO GERAL Na Tabela 6.3, são apresentadas as Estatísticas Básicas em relação ao componente da prova que avalia a Formação Geral dos estudantes concluintes.
A nota de Formação Geral é obtida levando -se em conta os acertos nas oito questões objetivas para o curso de Geografia ( Licenciatura ), já para o curso de Geografia ( Bacharelado ) foram levadas em conta sete questões objetivas ( uma foi anulada pela CAA ) e a nota Média obtida nas Questões estatisticamente significativa, ao nível de 95 \%, apenas entre as Médias das Universidades e dos CEFET/IF.
A Média das notas da prova, como um todo, dos alunos de Bacharelado, foi 43,2 ( Tabela 6.1 ).
Em relação a Categoria Administrativa, os estudantes de Bacharelado das IES Públicas obtiveram Média mais alta ( 43,4 ) que a Média nacional, e os das IES Privadas obtiveram Média mais baixa ( 40,9 ).
Observa -se que não diferença estatisticamente significativa entre as Médias das notas das IES Públicas e as das Privadas.
A diferença entre as Médias das regiões Sul e Norte ( 7,4 ), a maior e a menor Média, é superior à diferença entre a das IES Públicas e Privadas ( 2,5 ), caracterizando -se uma maior diversidade regional do que administrativa.
O Desvio padrão para as IES Públicas ( 13,1 ) foi inferior ao do Brasil, como um todo ( 13,2 ), indicando uma dispersão um pouco menor das notas nesta Categoria Administrativa.
No tocante a Organização Acadêmica, somente nas Universidades e nos Centros Universitários foram oferecidos cursos Bacharelado.
A Média das Universidades e dos Centros Universitários foram, respectivamente, 43,3 e 40,3.
Constata -se que não existe diferença estatisticamente significativa, ao nível de 95 \%, entre as Médias desses tipos de Organização Acadêmica. 6.1.2 ESTATÍSTICAS BÁSICAS NO COMPONENTE DE FORMAÇÃO GERAL Na Tabela 6.3, são apresentadas as Estatísticas Básicas em relação ao componente da prova que avalia a Formação Geral dos estudantes concluintes.
A nota de Formação Geral é obtida levando -se em conta os acertos nas oito questões objetivas para o curso de Geografia ( Licenciatura ), já para o curso de Geografia ( Bacharelado ) foram levadas em conta sete questões objetivas ( uma foi anulada pela CAA ) e a nota Média obtida nas Questões Discursivas 1 e 2, para ambas Habilitações, sendo a nota de cada questão discursiva obtida levando -se em conta 80 \% da nota relativa ao conteúdo e 20 \% da nota relativa ao desempenho linguístico.
Os estudantes de Licenciatura, de todo o Brasil, obtiveram desempenho médio de 34,6.
Quanto à variabilidade, o Desvio padrão das notas dos estudantes de Licenciatura do Brasil foi 15,2.
A maior Média foi obtida na região Sudeste ( 38,5 ) e a menor, na região Norte ( 31,4 ).
As demais Médias foram : 33,3 nas regiões Nordeste e Centro-Oeste, e 32,7 na região Sul.
Já o maior Desvio padrão foi obtido na região Sudeste ( 15,9 ), e o menor, na região Norte ( 14,1 ).
Os demais desvios foram : 14,5, na região Nordeste, 14,8, na região Sul e 14,9, na região Centro-Oeste.
A maior nota no Componente de Formação Geral da prova do Enade/2021 foi 84,9, obtida por, pelo menos, um estudante na região Sudeste.
A menor nota Máxima foi obtida na região Norte ( 73,4 ).
A Mediana do Brasil, como um todo, foi 34,2, sendo a menor Mediana encontrada na região Norte ( 30,1 ), e a maior, na região Sudeste ( 38,3 ).
A nota Mínima nessa parte foi zero em todas as regiões.
Considerando -se as notas, segundo Grande Região, observa -se que existe diferença estatisticamente significativa entre as Médias das notas do Componente de Formação Geral obtidas na região Sudeste e nas demais regiões.
Já os alunos de cursos Presenciais de Licenciatura obtiveram desempenho médio de 36,5.
Quanto à variabilidade, o Desvio padrão das notas do conjunto dos estudantes das Licenciaturas Presenciais foi 15,1.
A maior Média foi obtida na região Sudeste ( 41,5 ), e a menor, na região Norte ( 31,4 ).
As demais Médias foram : 34,4 na região Nordeste, 40,0 na região Sul e 33,3 na região Centro-Oeste.
O maior Desvio padrão foi obtido na região Sudeste ( 15,3 ), e o menor, nas regiões Norte ( 14,1 ).
Os demais desvios foram : 14,4 na região Nordeste, 14,7 na região Sul e 14,9 na região Centro-Oeste.
A nota Máxima no Componente de Formação Geral da prova de alunos de cursos Presenciais de Licenciatura foi 84,8, obtida por, pelo menos, um aluno na região Sul.
A menor nota Máxima foi obtida na região Norte ( 73,4 ).
A Mediana do conjunto de Licenciaturas Presenciais foi 36,5, sendo a menor Mediana encontrada na região Norte ( 30,1 ), e a maior, encontrada na região Sudeste ( 42,1 ).
A nota Mínima nesta parte foi zero em quase todas as regiões, com exceção da região Sul ( 2,0 ).
Considerando -se as notas segundo Grande Região, observa -se que não existe diferença estatisticamente significativa entre as Média da região Centro-Oeste e das regiões Norte e Nordeste, assim como entre as regiões Sudeste e Sul, nas notas das Licenciaturas Presenciais, no Componente de Formação Geral.
Apenas na região Norte não houve participação de estudantes de cursos a Distância de Licenciatura.
A Média das notas do Componente de Formação Geral dos alunos de cursos a Distância de Licenciatura, foi 32,0, mais baixa do que os alunos de cursos Presenciais ( 36,5 ) sendo que os alunos da região Nordeste obtiveram a Média mais baixa ( 27,0 ), e os da região Sudeste obtiveram a Média mais alta ( 35,1 ).
As demais Médias foram : 31,1 na região Sul e 33,3 na região Centro-Oeste.
O Desvio padrão das Licenciaturas a Distância foi 15,0, sendo o menor obtido na região Nordeste ( 13,5 ) e o maior na região Sudeste ( 15,9 ).
A maior Mediana foi a da região Sudeste ( 34,5 ) e a menor foi a da região Nordeste ( 25,8 ).
A maior nota Máxima foi 84,9, obtida na região Sudeste, registrada como a nota Máxima de todas as Licenciaturas.
As notas Mínimas foram zero em todas as regiões.
Na Modalidade de Educação a Distância de Licenciatura as diferenças entre Médias obtidas na região Nordeste e nas regiões Sudeste, Sul e Centro-Oeste são estatisticamente significativas, assim como entre as regiões Sudeste e Sul.
Os estudantes de Bacharelado, de todo o Brasil, obtiveram desempenho médio de 41,7 ( Tabela 6.3 ).
Quanto à variabilidade, o Desvio padrão das notas dos estudantes de Bacharelado do Brasil, como um todo, foi 15,1.
A maior Média foi obtida na região Sudeste ( 44,0 ), e a menor, na região Centro-Oeste ( 36,1 ).
As demais Médias foram : 38,0 na região Norte, 40,1, na região Nordeste e 42,4 na região Sul.
Já o maior Desvio padrão foi obtido na região Norte ( 16,6 ) e o menor, na região Nordeste ( 14,0 ).
Na região Sudeste o Desvio padrão foi 15,5, na região Sul foi 14,7 e na região Centro-Oeste foi 14,9.
A maior nota no Componente de Formação Geral da prova do Enade/2021 foi 86,2, obtida por, pelo menos, um estudante de Bacharelado da região Sudeste.
A menor nota Máxima foi obtida na região Centro-Oeste ( 71,5 ).
A Mediana do Brasil, como um todo, foi 42,0, sendo a menor Mediana encontrada na região Centro-Oeste ( 36,2 ), e a maior, na região Sudeste ( 44,7 ).
A nota Mínima nessa parte foi diferente de zero em três das cinco Grandes Regiões : Norte ( 1,0 ), Nordeste ( 3,2 ) e Sul ( 7,4 ).
Considerando -se as notas, segundo Grande Região, observa -se que existe diferença estatisticamente significativa entre as Médias das notas do Componente de Formação Geral obtidas na região Sudeste e nas regiões Norte, Nordeste e Centro-Oeste, como também entre as Médias da região Centro-Oeste e das regiões Nordeste e Sul.
A Média das notas da prova como um todo para os alunos de cursos Presenciais de Bacharelado, foi 42,1, sendo que os alunos da região Centro-Oeste obtiveram a Média mais baixa ( 36,1 ), e os da região Sudeste obtiveram a Média mais alta ( 44,6 ).
As demais Médias foram : 38,0 na região Norte, 40,1 na região Nordeste e 43,5 na região Sul.
A nota Máxima dos cursos Presenciais de Bacharelado foi obtida na região Sudeste ( 86,2 ) e a menor, na região Centro-Oeste ( 71,5 ).
A maior Mediana foi obtida na região Sudeste ( 45,0 ), e a menor, na Centro-Oeste ( 36,2 ).
A nota Mínima ( 0,0 ) foi obtida apenas na região Centro-Oeste.
Considerando -se as notas segundo Grande Região, não existe diferença estatisticamente significativa ao nível de 95 \% entre as Médias obtidas na região Norte e nas regiões Nordeste e Centro-Oeste, assim como entre as regiões Sudeste e Sul.
Somente nas regiões Sudeste e Sul foram oferecidos cursos a Distância de Bacharelado.
Respectivamente, para o Brasil, a região Sudeste e a região Sul as Médias foram 34,5, 31,6 e 36,0 ; o Desvio padrão foi 14,3, 15,5 e 13,5 ; a nota Mínima foi 0,0, 0,0 e 11,3 ; a Mediana foi 32,8, 29,1 e 36,0 ; e a nota Máxima foi 74,5, 64,5 e 74,5.
Não existe diferença estatisticamente significativa ao nível de 95 \% entre as Médias, obtidas nas regiões Sudeste e Sul.
O Gráfico 6.2 propicia a avaliação do desempenho dos estudantes do Componente de Formação Geral, a partir do histograma da distribuição das notas correspondentes.
A distribuição é unimodal para as notas de Licenciatura com moda em ( 30 ; 40 ], o mesmo intervalo modal da distribuição de notas da prova, como um todo ( Gráfico 6.1 ).
Para as notas do Bacharelado, a moda ocorre no intervalo ( 40 ; 50 ], o mesmo intervalo modal da distribuição de notas da prova, como um todo ( Gráfico 6.1 ).
Nota -se, ainda, que no Gráfico 6.2, as notas apresentam dispersão maior, quando comparadas a do Gráfico 6.1 ( distribuição das notas da prova ), confirmado pela comparação dos Desvios padrões que, em Formação Geral, aumenta de 14,3 para a nota da prova, como um todo ( Tabela 6.1 ), para 15,2 do Componente de Formação Geral ( Tabela 6.3 ), nas Licenciaturas e nos Bacharelados, aumenta de 13,2 para a nota da prova, como um todo ( Tabela 6.1 ), para 15,1 do Componente de Formação Geral ( Tabela 6.3 ).
Na Tabela 6.4, são apresentadas as informações referentes ao desempenho dos concluintes do Componente de Formação Geral, em diferentes agregações : Categoria Administrativa, Organização Acadêmica e Habilitação.
Primeiramente, com respeito aos alunos de Licenciatura, observa -se que existe diferença estatisticamente significativa entre as Médias dos tipos de Categoria Administrativa : a maior Média obtida por estudantes de Licenciatura de IES Públicas ( 35,5 ) e a menor, pelos de IES Privadas ( 32,7 ).
Nota -se que não há diferença estatisticamente significativa apenas entre as Médias das Universidades ( 34,9 ) e dos CEFET/IF ( 35,3 ).
Nas demais Organizações Acadêmicas, a Média foi : 33,2 para os Centros Universitários e 29,0 para as Faculdades.
Dentre os alunos de Bacharelado, observa -se que existe diferença estatisticamente significativa entre as Médias dos tipos de Categoria Administrativa : a maior Média obtida por estudantes de Bacharelado de IES Públicas ( 42,1 ) e a menor, pelos de IES Privadas ( 37,2 ).
No tocante a Organização Acadêmica, somente nas Universidades e nos Centros Universitários foram oferecidos cursos Bacharelado.
A Média das Universidades e dos Centros Universitários foram, respectivamente, 41,9 e 35,8.
Constata -se que existe diferença estatisticamente significativa, ao nível de 95 \%, entre as Médias desses tipos de Organização Acadêmica. 6.1.3 ESTATÍSTICAS BÁSICAS DO COMPONENTE DE CONHECIMENTO ESPECÍFICO Na Tabela 6.5, são apresentadas as Estatísticas Básicas referentes ao Componente de Conhecimento Específico da Área de Geografia ( Licenciatura/Bacharelado ) por Grande Região, segundo Habilitação e Modalidade de Ensino.
A nota do componente de Conhecimentos Específicos leva em conta as notas da parte objetiva e a média das notas obtidas nas três questões discursivas, parte discursiva da Área.
A Média do desempenho dos estudantes de Licenciatura do Brasil, como um todo, foi 42,9.
A maior Média foi obtida na região Sudeste ( 46,9 ), e a menor, na região Norte ( 37,5 ).
As demais Médias foram : 41,1 nas regiões Nordeste e Centro-Oeste, e 42,1 na região Sul.
Quanto à variabilidade das notas, o Desvio padrão do Brasil, como um todo para os alunos de Licenciatura, foi 16,0, sendo o maior Desvio padrão observado na região Sudeste ( 16,9 ) e o menor, na região Norte ( 13,7 ).
Os demais desvios foram : 15,0 na região Nordeste, 15,6 na região Sul e 15,9 na região Centro-Oeste.
A Mediana das notas dos estudantes de Licenciatura de todo o Brasil foi 42,8.
A maior Mediana ocorreu na região Sudeste ( 47,3 ), e a menor, na região Norte ( 36,0 ).
As demais Medianas foram : 42,0 nas regiões Nordeste e Sul, e 40,2 na região Centro-Oeste.
A nota Máxima do Brasil, como um todo, foi 97,5, sendo obtida por pelo menos um estudante de Licenciatura na região Sudeste.
As demais notas Máximas foram : 75,7 na região Norte, 89,3 na região Nordeste, 90,5 na região Sul e 83,5 na região Centro-Oeste.
A nota Mínima foi zero em quase todas as regiões, com exceção da região Norte ( 4,3 ).
Observa -se que não existe diferença estatisticamente significativa entre as regiões Nordeste, Sul e Centro-Oeste, na análise das Médias das notas do Componente de Conhecimento Específico.
Já para os alunos dos cursos Presenciais de Licenciatura, a Média do desempenho de Conhecimento Específico foi 44,5.
A maior Média foi obtida na região Sudeste ( 50,0 ), e a menor, na região Norte ( 37,5 ).
As demais Médias foram : 42,1 na região Nordeste, 49,5 na região Sul e 41,0 na região Centro-Oeste.
Quanto à variabilidade das notas, o Desvio padrão do conjunto das Licenciaturas Presenciais foi 16,1, sendo o maior Desvio padrão observado na região Sudeste ( 16,6 ), e o menor, na região Norte ( 13,7 ).
Os demais desvios foram : 14,9 na região Nordeste e 15,9 nas regiões Sul e Centro-Oeste.
A Mediana das notas dos estudantes de cursos Presenciais de Licenciatura de todo o Brasil foi 44,5.
A maior Mediana ocorreu na região Sudeste ( 51,8 ), e a menor, na região Norte ( 36,0 ).
As demais Medianas foram : 42,5 na região Nordeste, 50,5 na região Sul e 40,2 na região Centro-Oeste.
A nota Máxima das Licenciaturas Presenciais em Geografia foi 89,3, sendo obtida por, pelo menos, um aluno na região Nordeste.
As demais notas Máximas foram : 75,7 na região Norte, 88,7 na região Sudeste, 88,5 na região Sul e 83,5 na região Centro-Oeste.
A nota Mínima foi de zero nas regiões Nordeste e Sudeste, e nas demais regiões, foi 4,3.
Considerando -se as Licenciaturas oferecidas presencialmente, observa -se que não existe diferença estatisticamente significativa entre as regiões Nordeste e Centro-Oeste, na análise das Médias das notas do Componente de Conhecimento Específico, assim como entre as regiões Sudeste e Sul.
Já para os alunos de cursos Presenciais de Bacharelado, a Média do desempenho de Conhecimento Específico foi 43,9.
A maior Média foi obtida na região Sul ( 46,7 ), e a menor, Cursos a Distância de Licenciatura em Geografia só não foram ofertados na região Norte.
A Média do desempenho de Conhecimento Específico foi 40,9.
A maior Média foi obtida na região Sudeste ( 43,2 ), e a menor, na região Nordeste ( 35,5 ).
As demais Médias foram : 40,6 na região Sul e 41,3 na região Centro-Oeste.
Quanto à variabilidade das notas, o Desvio padrão do conjunto das Licenciaturas a Distância foi 15,6, sendo o maior Desvio padrão observado na região Sudeste ( 16,5 ), e o menor, na região Nordeste ( 14,3 ).
Os demais desvios foram : 15,1 na região Sul e 16,0 na região Centro-Oeste.
A Mediana das notas dos estudantes de cursos a Distância de Licenciatura de todo o Brasil foi 40,2.
A maior Mediana ocorreu na região Sudeste ( 42,5 ), e a menor, na região Nordeste ( 35,5 ).
As demais Medianas foram : 39,8 na região Sul e 40,0 na região Centro-Oeste.
A nota Máxima das Licenciaturas a Distância em Geografia foi 97,5, sendo obtida por, pelo menos, um aluno na região Sudeste.
As demais notas Máximas foram : 75,0 na região Nordeste, 90,5 na região Sul e 83,0 na região Centro-Oeste.
A nota Mínima foi de zero nas regiões Sudeste, Sul e Centro-Oeste, enquanto na região Nordeste foi 4,3.
A Média do desempenho dos estudantes de Bacharelado do Brasil, como um todo, foi 43,7.
A maior Média foi obtida na região Sul ( 46,0 ), e a menor, na região Norte ( 37,6 ).
As demais Médias foram : 41,9, na região Nordeste ; 45,2 na região Sudeste ; e 39,5, na região Centro-Oeste.
Quanto à variabilidade das notas, o Desvio padrão do Brasil, como um todo para os alunos de Bacharelado, foi 14,8, sendo o maior Desvio padrão observado na região Centro-Oeste ( 15,0 ) e o menor, na região Nordeste ( 14,0 ).
Os demais desvios foram : 14,9 na região Norte, 14,8 na região Sudeste e 14,5 na região Sul.
A Mediana das notas dos estudantes de Bacharelado de todo o Brasil foi 43,7.
A maior Mediana ocorreu nas regiões Sudeste e Sul ( 46,3 ), e a menor, na região Norte ( 37,8 ).
As demais Medianas foram : 42,2 na região Nordeste e 41,5 na região Centro-Oeste.
A nota Máxima do Brasil, como um todo, foi 83,7, sendo obtida por, pelo menos, um estudante de Bacharelado na região Sul.
As demais notas Máximas foram : 72,2, na região Norte ; 77,8, na região Nordeste ; 82,8, na região Sudeste e 75,7 na região Centro-Oeste.
A nota Mínima foi zero nas regiões Sudeste e Centro-Oeste, enquanto na região Norte foi 6,8, na região Nordeste foi 7,7 e na região Sul foi 1,0.
Considerando -se as Bacharelado oferecidas presencialmente, observa -se que não existe diferença estatisticamente significativa entre as regiões Norte, Nordeste e Centro-Oeste, na análise das Médias das notas do Componente de Conhecimento Específico, assim como entre as regiões Sudeste e Sul.
Já para os alunos de cursos Presenciais de Bacharelado, a Média do desempenho de Conhecimento Específico foi 43,9.
A maior Média foi obtida na região Sul ( 46,7 ), e a menor, na região Norte ( 37,6 ).
As demais Médias foram : 41,9 na região Nordeste, 45,6 na região Sudeste e 39,5 na região Centro-Oeste.
Quanto à variabilidade das notas, o Desvio padrão do conjunto das Bacharelado Presenciais foi 14,7, sendo o maior Desvio padrão observado na região Centro-Oeste ( 15,0 ), e o menor, na região Nordeste ( 14,0 ).
Os demais desvios foram : 14,9 na região Norte, 14,7 na região Sudeste e 14,3 na região Sul.
A Mediana das notas dos estudantes de cursos Presenciais de Bacharelado de todo o Brasil foi 44,1.
A maior Mediana ocorreu na região Sul ( 47,3 ), e a menor, na região Norte ( 37,8 ).
As demais Medianas foram : 42,2 na região Nordeste, 46,7 na região Sudeste e 41,5 na região Centro-Oeste.
A nota Máxima das Bacharelado Presenciais em Geografia foi 83,7, sendo obtida por, pelo menos, um aluno na região Sul.
As demais notas Máximas foram : 72,2 na região Norte, 77,8 na região Nordeste, 82,8 na região Sudeste e 75,7 na região Centro-Oeste.
A nota Mínima foi de zero nas regiões Sudeste e Centro-Oeste, e nas demais regiões foi : 6,8 na região Norte, 7,7 na região Nordeste e 1,0 na região Sul.
Somente nas regiões Sudeste e Sul foram oferecidos cursos a Distância de Bacharelado.
Respectivamente, para o Brasil, a região Sudeste e a região Sul as Médias foram 40,3, 37,5 e 41,9 ; o Desvio padrão foi 15,6, 15,8 e 15,4 ; a nota Mínima foi 5,1, 5,1 e 12,6 ; a Mediana foi 39,9, 38,7 e 41,2 ; e a nota Máxima foi 78,0, 78,0 e 77,3.
Não existe diferença estatisticamente significativa ao nível de 95 \% entre as Médias, obtidas nas regiões Sudeste e Sul.
Assim como nos Gráficos 6.1 e 6.2, no Gráfico 6.3, apresentado a seguir, é possível ser feita uma avaliação do desempenho de concluintes de Geografia, segundo a Habilitação, em relação ao Componente de Conhecimento Específico, com um histograma da distribuição das notas correspondentes.
Esta também é uma distribuição unimodal para ambas as Habilitações, e o grupo modal é o ( 40 ; 50 ], o mesmo do grupo modal da prova, como um todo ( Gráfico 6.1 ) e da Componente de Conhecimento Específico ( Gráfico 6.3 ) para os estudantes de Bacharelado.
Já para os estudantes de Licenciatura, corresponde ao intervalo imediatamente acima do grupo modal da prova, como um todo ( Gráfico 6.1 ) e da Componente de Conhecimento Específico ( Gráfico 6.3 ).
Na Tabela 6.6, é apresentada uma comparação dos resultados em relação a Categoria Administrativa, a Organização Acadêmica e segundo a Habilitação, agora levando -se em conta o desempenho de estudantes do Componente de Conhecimento Específico da prova.
No que se refere a Organização Acadêmica, para os alunos de Licenciatura, a maior Média foi a dos CEFET/IF ( 44,7 ), vindo a seguir a das Universidades ( 43,3 ), a dos Centros Universitários ( 41,1 ) e, por fim, a das Faculdades ( 36,5 ).
O maior Desvio padrão, e superior ao valor para o Brasil, como um todo, foi o dos Centros Universitários ( 16,1 ).
Os Centros Universitários obtiveram a maior nota Máxima ( 97,5 ).
As Universidades obtiveram nota Máxima 89,3, as Faculdades, 76,0, e os CEFET/IF, 87,5.
As Medianas foram : 43,0 nas Universidades, 40,5 nos Centros Universitários, 35,2 nas Faculdades e 45,5 nos CEFET/IF.
A nota Mínima foi zero para as Universidades e os Centros Universitários, enquanto foi 6,2 para as Faculdades e 5,0 para os CEFET/IF.
Observa -se que não existe diferença estatisticamente significativa, ao nível de 95 \%, no Componente de Conhecimento Específico apenas entre as notas das Universidades e dos CEFET/IF.
Quanto a Categoria Administrativa, observa -se um comportamento semelhante àquele da parte de Formação Geral e à prova, como um todo, ou seja, existe diferença estatisticamente significativa entre as Médias das IES Públicas ( 43,8 ) e as das IES Privadas ( 41,2 ).
No tocante a Organização Acadêmica, somente nas Universidades e nos Centros Universitários foram oferecidos cursos Bacharelado.
A Média das Universidades e dos Centros Universitários foram, respectivamente, 43,8 e 41,8.
Constata -se que não existe diferença estatisticamente significativa, ao nível de 95 \%, entre as Médias desses tipos de Organização Acadêmica.
O Desvio padrão foi 14,7 para as Universidades e 15,7 para os Centros Universitários.
As Universidades obtiveram a maior nota Máxima ( 83,7 ), enquanto os Centros Universitários obtiveram nota Máxima 77,3.
A Mediana foi 43,9 nas Universidades e 41,7 nos Centros Universitários.
A nota Mínima foi zero para as Universidades e 5,8 para os Centros Universitários.
Quanto a Categoria Administrativa, observa -se que não existe diferença estatisticamente significativa entre as Médias das IES Públicas ( 43,8 ) e as das IES Privadas ( 42,1 ). 6.2 ANÁLISE DAS QUESTÕES OBJETIVAS Esta seção apresenta estatísticas selecionadas e histogramas das Questões Objetivas de Formação Geral para Licenciatura ( 6.2.1 ) e para Bacharelado ( 6.2.2 ) e de Conhecimento Específico para Licenciatura ( 6.2.3 ) e para Bacharelado ( 6.2.4 ).
São, também, apresentadas e comparadas as médias das subpopulações caracterizadas por Grande Região. 21 6.2.1 COMPONENTE DE FORMAÇÃO GERAL – GEOGRAFIA ( LICENCIATURA ) Na Tabela 6.7a, são apresentadas as Estatísticas Básicas relativas às oito Questões Objetivas do Componente da prova que abrange a Formação Geral dos estudantes de Licenciatura.
A Média do Brasil foi 39,6.
A menor Média foi encontrada na região Norte ( 35,1 ), e a maior, na região Sudeste ( 43,8 ).
As demais Médias foram : 38,2 na região Nordeste, 37,9 na região Sul e 38,1 na região Centro-Oeste.
O Desvio padrão do Brasil foi 19,7, sendo o maior Desvio padrão encontrado na região Sudeste ( 20,2 ), e o menor Desvio padrão encontrado na região Norte ( 17,8 ).
Os demais desvios foram : 19,2 na região Nordeste, 19,3 na região Sul e 19,1 na região Centro-Oeste.
A Mediana ( 37,5 ), a nota Máxima ( 100,0 ) e a nota Mínima ( 0,0 ) foram iguais para todas as regiões, sendo a nota Máxima da região Norte ( 87,5 ) a única exceção.
Na Tabela 6.8a, são apresentados o Índice de Facilidade e o Índice de Discriminação ( Ponto-Bisserial ) para cada uma das Questões Objetivas do Componente de Formação Geral.
Quanto ao Índice de Facilidade, foram usadas as seguintes cores para diferenciar o nível de dificuldade da questão : Na Tabela 6.8a, são apresentados o Índice de Facilidade e o Índice de Discriminação ( Ponto-Bisserial ) para cada uma das Questões Objetivas do Componente de Formação Geral.
Quanto ao Índice de Facilidade, foram usadas as seguintes cores para diferenciar o nível de dificuldade da questão : Já quanto ao Índice de Discriminação, foram usadas as seguintes cores para qualificar a questão : As Questões Objetivas do Componente de Formação Geral, segundo o Índice de Facilidade, foram assim avaliadas : nenhuma das oito questões foi anulada pela Comissão Assessora de Área.
Nenhuma teve o Índice de Facilidade classificado como Muito fácil ou como Fácil.
Cinco questões foram consideradas com índice de dificuldade Médio, situando -se no intervalo entre 0,41 a 0,60 do Índice de Facilidade, ou seja, houve entre 41,0 \% e 60,0 \% de acertos, enquanto três questões foram classificadas na categoria Difícil, apresentando acertos na faixa entre 16 \% e 40 \%.
Nenhuma questão apresentou menos de 16 \% de acertos, razão pela qual não houve questão classificada como Muito difícil.
O Índice de Facilidade variou de 0,20 a 0,58, e o de Discriminação, de 0,24 a 0,49.
Cinco questões, quanto ao Índice de Discriminação, foram tidas como Muito bom, duas questões foram tidas como Bom e uma questão foi tida como Médio.
Na Tabela 6.9a, é apresentada a distribuição das questões, levando -se em conta, simultaneamente, a classificação dos dois índices.
Como já comentado, nenhuma questão foi anulada.
As cinco questões que tiveram Índice de Discriminação Muito bom figuraram em 21 As fórmulas utilizadas para o coeficiente de correlação ponto-bisserial e para o índice de facilidade estão apresentadas no capítulo 1.
Azul para as questões classificadas com índice Muito fácil ( \&gt ; =0,86 ), verde para as questões classificadas com índice Fácil ( 0,61 a 0,85 ), amarelo para as questões classificadas com Médio ( 0,41 a 0,60 ), vermelho para as questões classificadas com índice Difícil ( 0,16 a 0,40 ) e roxo para as questões classificadas com Muito difícil ( \&lt ; =0,15 ).
Já quanto ao Índice de Discriminação, foram usadas as seguintes cores para qualificar a questão : As questões classificadas com índice Fraco receberam a cor vermelha ( \&lt ; =0,19 ), as classificadas com Médio receberam a cor amarela ( 0,20 a 0,29 ), as classificadas com Bom receberam a cor verde ( 0,30 a 0,39 ) e as classificadas com Muito bom ( \&gt ; =0,40 ) receberam a cor azul.
As Questões Objetivas do Componente de Formação Geral, segundo o Índice de Facilidade, foram assim avaliadas : nenhuma das oito questões foi anulada pela Comissão Assessora de Área.
Nenhuma teve o Índice de Facilidade classificado como Muito fácil ou como Fácil.
Cinco questões foram consideradas com índice de dificuldade Médio, situando -se no intervalo entre 0,41 a 0,60 do Índice de Facilidade, ou seja, houve entre 41,0 \% e 60,0 \% de acertos, enquanto três questões foram classificadas na categoria Difícil, apresentando acertos na faixa entre 16 \% e 40 \%.
Nenhuma questão apresentou menos de 16 \% de acertos, razão pela qual não houve questão classificada como Muito difícil.
O Índice de Facilidade variou de 0,20 a 0,58, e o de Discriminação, de 0,24 a 0,49.
Cinco questões, quanto ao Índice de Discriminação, foram tidas como Muito bom, duas questões foram tidas como Bom e uma questão foi tida como Médio.
Na Tabela 6.9a, é apresentada a distribuição das questões, levando -se em conta, simultaneamente, a classificação dos dois índices.
Como já comentado, nenhuma questão foi anulada.
As cinco questões que tiveram Índice de Discriminação Muito bom figuraram em apenas um dos níveis de dificuldade, todas classificadas na categoria Médio ( Questões 1, 3, 5, 6 e 7 ).
Em particular, as Questões 1 e 7 foram as que apresentaram o maior poder discriminatório, ambas com índice 0,49, e foram consideradas Médio em termos de Facilidade, com proporções de 0,43 e 0,58 acertos, respectivamente.
O máximo de acertos foi alcançado pelas Questões 6 e 7, ambas com um Índice de Facilidade de 0,58.
A Questão 8 foi considerada Difícil, com Índice de Facilidade 0,20 e Índice de Discriminação Médio ( 0,24 ).
No Gráfico 6.4a, para exemplificar, analisa -se o comportamento da Questão de número 3 de Formação Geral.
Trata -se de uma questão considerada com nível Médio de dificuldade, em relação à Facilidade, e a que obteve o quarto maior Índice de Discriminação dessa parte da prova.
Neste gráfico, cada uma das cinco curvas representa o percentual de respostas em determinada alternativa da questão, em função do número de acertos dos estudantes nessa parte da prova ( Formação Geral/Múltipla Escolha ), antes de possíveis eliminações pelo critério Ponto-Bisserial.
Como já foi mencionado, nenhuma Questão Objetiva do Componente de Formação Geral foi anulada pela Comissão Assessora de Área, por isso o número de acertos alcança no máximo oito para esse componente.
A curva em verde corresponde à alternativa B, a alternativa correta para essa questão.
Observa -se que entre os estudantes com menor número de acertos, nessa parte do exame, a situação mais frequente foi a escolha da alternativa E ( em vermelho ), incorreta.
Por exemplo, entre os estudantes que acertaram duas respostas, 46,3 \% escolheram a alternativa E ( em vermelho ), 27,0 \% escolheram a alternativa B ( em verde, gabarito ), 12,1 \% escolheram a alternativa D ( em laranja ), 9,2 \% escolheram a alternativa A ( em azul ) e 5,0 \% escolheram a alternativa C ( em preto ).
Dentre os que acertaram duas respostas entre as questões de múltipla escolha de Formação Geral, 0,1 \% deixou a questão em branco e 0,2 \% marcou mais de uma alternativa, invalidando a questão.
À medida que o número de acertos aumenta, indicando desempenho melhor nessa parte da prova, aumenta, concomitantemente, a proporção de estudantes que selecionaram a alternativa correta B, atingindo 100 \% para os estudantes com oito acertos.
Essa análise permite verificar como a questão discriminou os grupos de desempenho, justificando -se o alto índice ( 0,46 ) obtido na questão.
Cumpre notar que não é possível inferir deste gráfico nem o Índice de Facilidade, que seria uma média da proporção ponderada pela quantidade de estudantes com cada uma das notas, nem o Índice de Discriminação Ponto-Bisserial, por razão equivalente.
No caso extremo, no qual a grande concentração dos acertos dos estudantes fosse abaixo de quatro, o Índice de Facilidade seria obrigatoriamente abaixo de 60,0 \% ( neste exemplo ).
Caso a concentração fosse em seis acertos ou mais, o índice seria obrigatoriamente acima de 80,0 \%.
Os gráficos relativos às demais questões de Formação Geral constam do Anexo IA. 6.2.2 COMPONENTE DE FORMAÇÃO GERAL – GEOGRAFIA ( BACHARELADO ) Como nenhuma questão foi anulada pela Comissão Assessora de Área, na Tabela 6.7b, são apresentadas as Estatísticas Básicas relativas às oito Questões Objetivas válidas do Componente da prova que abrange a Formação Geral dos estudantes de Bacharelado.
A Média do Brasil foi 47,6.
A menor Média foi encontrada na região Centro-Oeste ( 43,0 ), e a maior, na região Sudeste ( 49,4 ).
As demais Médias foram : 46,3, na região Norte, 45,4, na região Nordeste e 49,2, na região Sul.
O Desvio padrão do Brasil foi 19,5, sendo o maior Desvio padrão encontrado na região Norte ( 21,3 ), e o menor, na região Nordeste ( 18,9 ).
Os demais desvios foram : 19,6 nas regiões Sudeste e Centro-Oeste, e 19,0 na região Sul.
A Mediana ( 50,0 ), a nota Máxima ( 100,0 ) e a nota Mínima ( 0,0 ) foram iguais para todas as regiões, sendo a nota Máxima e a Mediana da região Norte ( 87,5 e 37,5, respectivamente ) as únicas exceções.
Na Tabela 6.8b, são apresentados o Índice de Facilidade e o Índice de Discriminação ( Ponto-Bisserial ) para cada uma das Questões Objetivas do Componente de Formação Geral.
A convenção das cores é a mesma da Tabela 6.8a.
As Questões Objetivas do Componente de Formação Geral, segundo o Índice de Facilidade, foram assim avaliadas : das oito questões válidas, nenhuma teve o Índice de Facilidade classificado como Muito fácil, e três questões foram tidas como Fácil, por terem índice de acertos entre 0,61 a 0,85.
Duas questões foram consideradas com índice de dificuldade Médio, situando -se no intervalo entre 0,41 e 0,60 do Índice de Facilidade, ou seja, houve entre 41,0 \% e 60,0 \% de acertos, enquanto outras três questões foram classificadas na categoria Difícil, situando -se no intervalo entre 0,16 e 0,40.
Nenhuma questão apresentou menos de 16 \% de acertos, razão pela qual não houve questão classificada como Muito difícil.
O Índice de Facilidade variou de 0,19 a 0,68 e o de Discriminação, de 0,30 a 0,48.
Cinco das oito questões válidas, quanto ao Índice de Discriminação, foram tidas como Muito bom e três como Bom.
Desvio padrão encontrado na região Norte ( 21,3 ), e o menor, na região Nordeste ( 18,9 ).
Os demais desvios foram : 19,6 nas regiões Sudeste e Centro-Oeste, e 19,0 na região Sul.
A Mediana ( 50,0 ), a nota Máxima ( 100,0 ) e a nota Mínima ( 0,0 ) foram iguais para todas as regiões, sendo a nota Máxima e a Mediana da região Norte ( 87,5 e 37,5, respectivamente ) as únicas exceções.
Na Tabela 6.8b, são apresentados o Índice de Facilidade e o Índice de Discriminação ( Ponto-Bisserial ) para cada uma das Questões Objetivas do Componente de Formação Geral.
A convenção das cores é a mesma da Tabela 6.8a.
As Questões Objetivas do Componente de Formação Geral, segundo o Índice de Facilidade, foram assim avaliadas : das oito questões válidas, nenhuma teve o Índice de Facilidade classificado como Muito fácil, e três questões foram tidas como Fácil, por terem índice de acertos entre 0,61 a 0,85.
Duas questões foram consideradas com índice de dificuldade Médio, situando -se no intervalo entre 0,41 e 0,60 do Índice de Facilidade, ou seja, houve entre 41,0 \% e 60,0 \% de acertos, enquanto outras três questões foram classificadas na categoria Difícil, situando -se no intervalo entre 0,16 e 0,40.
Nenhuma questão apresentou menos de 16 \% de acertos, razão pela qual não houve questão classificada como Muito difícil.
O Índice de Facilidade variou de 0,19 a 0,68 e o de Discriminação, de 0,30 a 0,48.
Cinco das oito questões válidas, quanto ao Índice de Discriminação, foram tidas como Muito bom e três como Bom.
Na Tabela 6.9b, é apresentada a distribuição das questões, levando -se em conta, simultaneamente, a classificação dos dois índices.
As cinco questões que tiveram Índice de Discriminação Muito bom figuraram entre dois níveis de dificuldade Fácil e Médio : três classificadas na categoria Fácil ( Questões 1, 6 e 7 ) do Índice de Facilidade e duas, na categoria Médio ( Questões 3 e 5 ).
Em particular, a Questão 1 foi a que apresentou o maior poder discriminatório, com índice 0,48, e foi considerada Fácil em termos de Facilidade, com uma proporção de 0,61 acertos.
O máximo de acertos foi alcançado pela Questão 6, com um Índice de Facilidade de 0,68.
A Questão 8 foi considerada Difícil, com Índice de Facilidade 0,19.
Seu Índice de Discriminação foi o mais baixo dentre as oito questões válidas : Bom ( 0,30 ).
No Gráfico 6.4b, para exemplificar, analisa -se o comportamento da Questão de número 3 de Formação Geral.
Trata -se de uma questão considerada Médio, em relação à Facilidade e a que obteve o segundo maior Índice de Discriminação dessa parte da prova.
Neste gráfico, cada uma das cinco curvas representa o percentual de respostas em determinada alternativa da questão, em função do número de acertos dos estudantes nessa parte da prova ( Formação Geral/Múltipla Escolha ), antes de possíveis eliminações pelo Na Tabela 6.8b, são apresentados o Índice de Facilidade e o Índice de Discriminação ( Ponto-Bisserial ) para cada uma das Questões Objetivas do Componente de Formação Geral.
A convenção das cores é a mesma da Tabela 6.8a.
As Questões Objetivas do Componente de Formação Geral, segundo o Índice de Facilidade, foram assim avaliadas : das oito questões válidas, nenhuma teve o Índice de Facilidade classificado como Muito fácil, e três questões foram tidas como Fácil, por terem índice de acertos entre 0,61 a 0,85.
Duas questões foram consideradas com índice de dificuldade Médio, situando -se no intervalo entre 0,41 e 0,60 do Índice de Facilidade, ou seja, houve entre 41,0 \% e 60,0 \% de acertos, enquanto outras três questões foram classificadas na categoria Difícil, situando -se no intervalo entre 0,16 e 0,40.
Nenhuma questão apresentou menos de 16 \% de acertos, razão pela qual não houve questão classificada como Muito difícil.
O Índice de Facilidade variou de 0,19 a 0,68 e o de Discriminação, de 0,30 a 0,48.
Cinco das oito questões válidas, quanto ao Índice de Discriminação, foram tidas como Muito bom e três como Bom.
Na Tabela 6.9b, é apresentada a distribuição das questões, levando -se em conta, simultaneamente, a classificação dos dois índices.
As cinco questões que tiveram Índice de Discriminação Muito bom figuraram entre dois níveis de dificuldade Fácil e Médio : três classificadas na categoria Fácil ( Questões 1, 6 e 7 ) do Índice de Facilidade e duas, na categoria Médio ( Questões 3 e 5 ).
Em particular, a Questão 1 foi a que apresentou o maior poder discriminatório, com índice 0,48, e foi considerada Fácil em termos de Facilidade, com uma proporção de 0,61 acertos.
O máximo de acertos foi alcançado pela Questão 6, com um Índice de Facilidade de 0,68.
A Questão 8 foi considerada Difícil, com Índice de Facilidade 0,19.
Seu Índice de Discriminação foi o mais baixo dentre as oito questões válidas : Bom ( 0,30 ).
No Gráfico 6.4b, para exemplificar, analisa -se o comportamento da Questão de número 3 de Formação Geral.
Trata -se de uma questão considerada Médio, em relação à Facilidade e a que obteve o segundo maior Índice de Discriminação dessa parte da prova.
Neste gráfico, cada uma das cinco curvas representa o percentual de respostas em determinada alternativa da questão, em função do número de acertos dos estudantes nessa parte da prova ( Formação Geral/Múltipla Escolha ), antes de possíveis eliminações pelo critério Ponto-Bisserial.
A curva em verde corresponde à alternativa B, a alternativa correta para essa questão.
Observa -se que, entre os estudantes com menor número de acertos, nessa parte do exame, a situação mais frequente foi a escolha da alternativa E ( em vermelho ), incorreta.
Por exemplo, entre os estudantes que acertaram duas questões, 46,1 \% escolheram a alternativa E ( em vermelho ), 25,0 \% escolheram a alternativa B ( em verde, gabarito ), 13,9 \% escolheram a alternativa A ( em azul ), 8,2 \% escolheram a alternativa D ( em laranja ) e 6,4 \%, a C ( em preto ).
Dentre os que acertaram duas respostas entre as questões de múltipla escolha de Formação Geral, 0,4 \% deixou a questão em branco e 0,0 \% marcou mais de uma alternativa, invalidando a questão.
À medida que o número de acertos aumenta, indicando desempenho melhor nessa parte da prova, aumenta, concomitantemente, a proporção de estudantes que selecionaram a alternativa correta B, atingindo 100 \% para os estudantes com oito acertos.
Essa análise permite verificar como a questão discriminou os grupos de desempenho, justificando -se o alto índice ( 0,47 ) obtido na questão.
Cumpre notar que não é possível inferir deste gráfico nem o Índice de Facilidade, que seria uma média da proporção ponderada pela quantidade de estudantes com cada uma das notas, nem o Índice de Discriminação Ponto-Bisserial, por razão equivalente.
No caso extremo, no qual a grande concentração dos acertos dos estudantes fosse abaixo de quatro, o Índice de Facilidade seria obrigatoriamente abaixo de 60,0 \% ( neste exemplo ).
Caso a concentração fosse em seis acertos ou mais, o índice seria obrigatoriamente acima de 80,0 \%.
Os gráficos relativos às demais questões de Formação Geral constam do Anexo IB. 6.2.3 COMPONENTE DE CONHECIMENTO ESPECÍFICO – GEOGRAFIA ( LICENCIATURA ) Na Tabela 6.10a, são apresentadas as Estatísticas Básicas em relação às questões objetivas do Componente de Conhecimento Específico da prova de Geografia por Grande Região.
Nessa parte da prova, como será discutido ainda nesta seção, nenhuma das questões objetivas do Componente de Conhecimentos Específicos foi anulada pela CAA e sete questões não foram usadas no cômputo das notas por terem sido descartadas pelo critério do Ponto-Bisserial.
Assim, as notas foram calculadas com base em 20 das 27 Questões Objetivas válidas de Conhecimentos Específicos.
A Média do Brasil deste componente foi 47,3.
A menor Média foi observada na região Norte ( 41,1 ), e a maior, na região Sudeste ( 51,2 ).
O Desvio padrão de todo o Brasil foi 17,5, sendo o menor Desvio padrão encontrado na região Norte ( 15,0 ), e o maior, na região Sudeste ( 18,4 ).
A Mediana de todo o Brasil foi 45,0, o mesmo valor da Mediana encontrada nas regiões Nordeste, Sul e Centro-Oeste.
Na região Norte ( 40,0 ), o valor foi menor, e, na região Sudeste ( 55,0 ), o valor foi maior.
A nota Máxima da prova foi 100,00, obtida nas Questões Objetivas do Componente de Conhecimento Específico, por, pelo menos, um estudante da região Sudeste.
Nas demais regiões os valores foram : 85,0 na região Norte, 95,0 nas regiões Nordeste e Sul, e 90,0 na região Centro-Oeste.
A nota Mínima foi zero, obtida por pelo menos um estudante em quase todas as regiões, com exceção da região Norte ( 5,0 ).
Na Tabela 6.11a, são apresentados os Índices de Facilidade e de Discriminação ( Ponto-Bisserial ) das Questões Objetivas do Componente de Conhecimento Específico, para os estudantes de Geografia de Licenciatura.
Para facilitar a diferenciação das questões, serão usadas as mesmas cores da Tabela 6.8a e 6.8b para as diferentes classificações dos Índices de Facilidade e de Discriminação.
Observa -se que nenhuma questão foi anulada pela CAA, como já mencionado.
Sendo assim, as análises a seguir dizem respeito às 27 questões válidas.
A partir do Índice de Facilidade obtido, pode -se concluir que quase metade das questões objetivas da prova foram consideradas Difícil ou Muito difícil : das 27 questões válidas, 11 foram classificadas como Difícil e uma como Muito difícil.
Três questões foram classificadas como Fácil, e outras 12 consideradas como Médio.
Nenhuma questão foi classificada como Muito fácil.
Já quanto aos Índices de Discriminação das Questões Objetivas do Componente de Conhecimento Específico da prova, obtém -se como resultado a seguinte classificação : seis das 27 questões válidas foram consideradas boas, enquanto nove delas tiveram Índice de Discriminação Muito bom.
Assim, para 15 em 27 questões, os Índices de Discriminação foram Bom ou Muito bom.
Dentre as demais, cinco delas foram classificadas como Médio, e outras sete, como Fraco, sendo 12, por conseguinte, a quantidade de questões nos dois patamares mais baixos de Discriminação.
Constata -se, assim, que a prova – no que se refere ao Componente de Conhecimento Específico – possuía média capacidade de discriminar entre aqueles que dominam ou não o conteúdo.
O Índice de Facilidade variou de 0,11 a 0,78 e o de Discriminação, de -0,03 a 0,48.
Na Tabela 6.12a, é apresentada a distribuição das questões levando -se em conta, simultaneamente, a classificação dos dois índices.
Dentre as questões que alcançaram os maiores Índices de Discriminação, uma delas foi classificada com Índice de Discriminação Muito bom, a de número 27.
Seu índice foi 0,48, e, quanto ao Índice de Facilidade, foi classificada como Médio, com 57 \% dos estudantes marcando a opção correta.
A Questão de número 16 foi a mais difícil entre as 27 questões específicas válidas, com baixo Índice de Facilidade, apenas 11 \% de acertos.
Essa questão apresentou poder discriminatório muito baixo, -0,01, o que comprova ter sido a mais difícil para os estudantes.
Destaca -se, também, a Questão 22, com Índice de Facilidade 0,20, o que, em termos percentuais, corresponde a 20 \% de estudantes que responderam acertadamente, sendo, também, 0,08 seu Índice de Discriminação.
Tais questões foram, portanto, pelo critério Ponto-Bisserial, consideradas inadequadas.
Por isso, as Questões 16 e 22, além das Questões 9, 12, 17, 24 e 29, foram eliminadas do cômputo da nota final.
A título de exemplo das análises do comportamento das questões objetivas, no Gráfico 6.5a, analisa -se a Questão 27 do Componente de Conhecimento Específico.
Essa questão foi considerada, pelas respostas dos estudantes avaliados na prova, como uma questão com nível de Facilidade Médio, 0,57, ou seja, 57 \% dos estudantes assinalaram, acertadamente, a opção E correspondente ao gabarito.
Como já comentado, seu Índice de Discriminação foi igual a 0,48, classificado como Muito bom, o maior valor de Discriminação.
Neste gráfico, cada uma das cinco curvas representa o percentual de respostas em determinada alternativa da Questão 27, em função do número de acertos dos estudantes nessa parte da prova, antes de possíveis eliminações de questões pelo critério Ponto-Bisserial.
A alternativa correta E, representada no gráfico pela curva em vermelho, foi escolhida em maiores proporções pelos estudantes com desempenho melhor nessa parte da prova.
Já as alternativas incorretas, também denominadas distratores, foram selecionadas, principalmente, por aqueles com notas mais baixas.
No caso de Geografia ( Licenciatura ), como se observa no eixo horizontal do Gráfico 6.5a, nenhum estudante acertou mais do que 24 questões dentre as 27 questões objetivas válidas, considerando -se o gabarito original.
A eliminação das sete questões com Índice de Discriminação Fraco do cômputo da nota é que possibilitou a existência de nota Máxima igual a 100,0 por pelo menos um estudante da região Sudeste ( visto na Tabela 6.10a ).
Observa -se que a soma não é 100 \%, por causa das questões não respondidas ou com mais de uma opção marcada.
Aqueles com nota zero, 67,5 \% marcaram a opção D ( linha baixo, -0,01, o que comprova ter sido a mais difícil para os estudantes.
Destaca -se, também, a Questão 22, com Índice de Facilidade 0,20, o que, em termos percentuais, corresponde a 20 \% de estudantes que responderam acertadamente, sendo, também, 0,08 seu Índice de Discriminação.
Tais questões foram, portanto, pelo critério Ponto-Bisserial, consideradas inadequadas.
Por isso, as Questões 16 e 22, além das Questões 9, 12, 17, 24 e 29, foram eliminadas do cômputo da nota final.
A título de exemplo das análises do comportamento das questões objetivas, no Gráfico 6.5a, analisa -se a Questão 27 do Componente de Conhecimento Específico.
Essa questão foi considerada, pelas respostas dos estudantes avaliados na prova, como uma questão com nível de Facilidade Médio, 0,57, ou seja, 57 \% dos estudantes assinalaram, acertadamente, a opção E correspondente ao gabarito.
Como já comentado, seu Índice de Discriminação foi igual a 0,48, classificado como Muito bom, o maior valor de Discriminação.
Neste gráfico, cada uma das cinco curvas representa o percentual de respostas em determinada alternativa da Questão 27, em função do número de acertos dos estudantes nessa parte da prova, antes de possíveis eliminações de questões pelo critério Ponto-Bisserial.
A alternativa correta E, representada no gráfico pela curva em vermelho, foi escolhida em maiores proporções pelos estudantes com desempenho melhor nessa parte da prova.
Já as alternativas incorretas, também denominadas distratores, foram selecionadas, principalmente, por aqueles com notas mais baixas.
No caso de Geografia ( Licenciatura ), como se observa no eixo horizontal do Gráfico 6.5a, nenhum estudante acertou mais do que 24 questões dentre as 27 questões objetivas válidas, considerando -se o gabarito original.
A eliminação das sete questões com Índice de Discriminação Fraco do cômputo da nota é que possibilitou a existência de nota Máxima igual a 100,0 por pelo menos um estudante da região Sudeste ( visto na Tabela 6.10a ).
Observa -se que a soma não é 100 \%, por causa das questões não respondidas ou com mais de uma opção marcada.
Aqueles com nota zero, 67,5 \% marcaram a opção D ( linha laranja ) e 32,5 \% a opção A ( linha azul ).
A proporção de estudantes que selecionaram a resposta correta E aumenta gradativamente, chegando a atingir 100 \% para 21 acertos ou mais, enquanto a proporção dos que escolheram alternativas incorretas decai, a partir do primeiro acerto, em função do número de acertos nessa parte da prova.
Os gráficos relativos às demais questões do Conhecimento Específico constam do Anexo IA. 6.2.4 COMPONENTE DE CONHECIMENTO ESPECÍFICO – GEOGRAFIA ( BACHARELADO ) Na Tabela 6.10b, são apresentadas as Estatísticas Básicas em relação às questões objetivas do Componente de Conhecimento Específico da prova de Geografia por Grande Região.
Nessa parte da prova, nenhuma questão foi anulada pela Comissão Assessora da Área.
Além disso, como será discutido ainda nesta seção, cinco questões objetivas do Componente de Conhecimentos Específicos não foram usadas no cômputo das notas por terem sido descartadas pelo critério do Ponto-Bisserial.
Assim, as notas foram calculadas com base em 22 das 27 Questões Objetivas válidas de Conhecimentos Específicos.
A Média do Brasil deste componente foi 47,9.
A menor Média foi observada na região Norte ( 41,1 ), e a maior, na região Sul ( 50,5 ).
O Desvio padrão de todo o Brasil foi 16,2, sendo o menor Desvio padrão encontrado na região Nordeste ( 15,3 ), e o maior, na região Centro-Oeste ( 16,4 ).
A Mediana de todo o Brasil foi 50,0, o mesmo valor da Mediana encontrada nas regiões Sudeste e Sul.
Na região Norte a Mediana foi 40,9 e nas regiões Nordeste e Centro-Oeste, foi 45,5.
A nota Máxima da prova foi 95,5, obtida nas Questões Objetivas do Componente de Conhecimento Específico, por, pelo menos, um estudante da região Sul.
A nota Mínima foi zero nas regiões Sudeste, Sul e Centro-Oeste, enquanto na região Norte foi 4,5 e na Nordeste, foi 9,1.
Na Tabela 6.11b, são apresentados os Índices de Facilidade e de Discriminação ( Ponto-Bisserial ) das Questões Objetivas do Componente de Conhecimento Específico, para os estudantes de Geografia de Bacharelado.
Para facilitar a diferenciação das questões, serão usadas as mesmas cores da Tabela 6.8a, 6.8b e 6.11a para as diferentes classificações dos Índices de Facilidade e de Discriminação.
Observa -se que nenhuma questão foi anulada pela Comissão Assessora da Área ( CAA ) e que as análises dizem respeito às 27 questões válidas.
A partir do Índice de Facilidade obtido, pode -se concluir que pouco menos da metade das questões objetivas da prova foram considerados Difícil ou Muito difícil, das 27 questões válidas, dez questões foram classificadas como Difícil e uma como Muito difícil.
Duas questões foram classificadas como Fácil, e outras 14 consideradas como Médio.
Nenhuma questão foi classificada como Muito fácil. o menor Desvio padrão encontrado na região Nordeste ( 15,3 ), e o maior, na região Centro-Oeste ( 16,4 ).
A Mediana de todo o Brasil foi 50,0, o mesmo valor da Mediana encontrada nas regiões Sudeste e Sul.
Na região Norte a Mediana foi 40,9 e nas regiões Nordeste e Centro-Oeste, foi 45,5.
A nota Máxima da prova foi 95,5, obtida nas Questões Objetivas do Componente de Conhecimento Específico, por, pelo menos, um estudante da região Sul.
A nota Mínima foi zero nas regiões Sudeste, Sul e Centro-Oeste, enquanto na região Norte foi 4,5 e na Nordeste, foi 9,1.
Na Tabela 6.11b, são apresentados os Índices de Facilidade e de Discriminação ( Ponto-Bisserial ) das Questões Objetivas do Componente de Conhecimento Específico, para os estudantes de Geografia de Bacharelado.
Para facilitar a diferenciação das questões, serão usadas as mesmas cores da Tabela 6.8a, 6.8b e 6.11a para as diferentes classificações dos Índices de Facilidade e de Discriminação.
Observa -se que nenhuma questão foi anulada pela Comissão Assessora da Área ( CAA ) e que as análises dizem respeito às 27 questões válidas.
A partir do Índice de Facilidade obtido, pode -se concluir que pouco menos da metade das questões objetivas da prova foram considerados Difícil ou Muito difícil, das 27 questões válidas, dez questões foram classificadas como Difícil e uma como Muito difícil.
Duas questões foram classificadas como Fácil, e outras 14 consideradas como Médio.
Nenhuma questão foi classificada como Muito fácil.
Já quanto aos Índices de Discriminação das Questões Objetivas do Componente de Conhecimento Específico da prova, obtém -se como resultado a seguinte classificação : nove das 27 questões foram consideradas como Bom, enquanto cinco delas tiveram Índice de Discriminação Muito bom.
Assim, para 14 em 27 questões, os Índices de Discriminação foram Bom ou Muito bom.
Dentre as demais, oito delas foram classificadas como Médio, e cinco, como Fraco, sendo 13, por conseguinte, a quantidade de questões nos dois patamares mais baixos de discriminação.
Constata -se, assim, que a prova – no que se refere ao Componente de Conhecimento Específico – possuía média capacidade de discriminar entre aqueles que dominam ou não o conteúdo.
O Índice de Facilidade variou de 0,14 a 0,73 e o de Discriminação, de -0,03 a 0,47.
Na Tabela 6.12b, é apresentada a distribuição das questões levando -se em conta, simultaneamente, a classificação dos dois índices.
Dentre as questões que alcançaram os maiores Índices de Discriminação, uma delas obteve o maior deles, a de número 32.
Seu índice foi 0,47, e, quanto ao Índice de Facilidade, foi classificada como Médio, com 50 \% dos estudantes marcando a opção correta.
A Questão de número 16 foi a mais difícil entre as 27 questões específicas válidas, com baixo Índice de Facilidade, apenas 14 \% de acertos.
Essa questão apresentou poder discriminatório Médio, 0,22, o que comprova ter sido uma das mais difíceis para os estudantes.
Destacam -se, também, as Questões 9 e 24, com Índice de Facilidade 0,23, o que, em termos percentuais, corresponde a 23 \% de estudantes que responderam acertadamente.
Já -0,03 e 0,01 foram, respectivamente, os seus Índices de Discriminação.
Tais questões foram, portanto, pelo critério Ponto-Bisserial, consideradas inadequadas.
Por isso, as Questões 9 e 24, além das Questões 12, 14 e 17, foram eliminadas do cômputo da nota final.
A título de exemplo das análises do comportamento das questões objetivas, no Gráfico 6.5b, analisa -se a Questão 32 do Componente de Conhecimento Específico.
Essa questão foi considerada, pelas respostas dos estudantes avaliados na prova, como uma questão com nível de Facilidade Médio, 0,50, ou seja, 50 \% dos estudantes assinalaram, acertadamente, a opção E ( em vermelho ), correspondente ao gabarito.
Como já comentado, seu Índice de Discriminação foi igual a 0,47, classificado como Muito bom, o maior valor de discriminação.
Neste gráfico, cada uma das cinco curvas representa o percentual de respostas em determinada alternativa da Questão 32, em função do número de acertos dos estudantes nessa parte da prova, antes de possíveis eliminações de questões pelo critério Ponto-Bisserial.
A alternativa correta E, representada no gráfico pela curva em vermelho, foi escolhida em maiores proporções pelos estudantes com desempenho melhor nessa parte da prova.
Já as alternativas incorretas, também denominadas distratores, foram selecionadas, principalmente, por aqueles com notas mais baixas.
No caso de Geografia ( Bacharelado ), como se observa no eixo horizontal do Gráfico 6.5b, nenhum estudante acertou mais do que 23 questões dentre as 27 questões objetivas válidas, considerando -se o gabarito original.
Mesmo com a eliminação das cinco questões com Índice de Discriminação Fraco do cômputo da nota, nenhum participante atingiu a nota Máxima, igual a 100,0, nessa parte da prova ( visto na Tabela 6.10b ).
Observa -se que a soma não é 100 \%, por causa das questões não respondidas ou com mais de uma opção marcada.
Aqueles com um acerto, 50 \% marcaram a opção D ( em laranja ) e 25 \% a opção C ( em preto ).
A proporção de estudantes que selecionaram a resposta correta E aumenta gradativamente, chegando a atingir 100 \% para 22 acertos ou mais, enquanto a proporção dos que escolheram alternativas incorretas decai, a partir de três acertos, em função do número de acertos nessa parte da prova.
Os gráficos relativos às demais questões do Conhecimento Específico constam do Anexo IB. 6.3 ANÁLISE DAS QUESTÕES DISCURSIVAS Esta seção apresenta estatísticas selecionadas e histogramas das Questões Discursivas de Formação Geral ( 6.3.1 ) e Conhecimento Específico ( 6.3.2 ).
São, também, apresentadas e comparadas as médias de subpopulações, caracterizadas por Grande Região.
Cumpre notar que nem todas as questões passam pelo crivo dos corretores.
As respostas depois de digitalizadas passam por um pré-processamento de identificação automático.
Algumas destas questões são definidas como um “ branco automático ” ( que não são enviadas para os professores para correção, a não ser que na segunda fase, a de inspeção visual, seja constatada uma resposta ).
O critério usado neste pré-processamento, é verificar para cada questão, a partir da quantidade de PIXELS, o que provavelmente, é uma resposta em branco.
A partir da separação destes documentos que foram considerados em branco, eles são enviados para uma verificação visual, por uma equipe de colaboradores, devidamente treinados.
Deste modo, quaisquer erros, deste pré-processamento, são encontrados, e as questões preenchidas, são enviadas para a correção dos professores.
Eventualmente, algumas respostas em branco, não são detectadas neste pré-processamento automático.
A presença de pequenos riscos, ou manchas, podem fazer o programa não identificar uma questão como um possível branco.
Estas são enviadas para a correção, mas são classificadas como “ branco ”, pelos professores quando da correção.
No caso da Área de Geografia ( Licenciatura ), foram identificados como brancos “ automáticos ”, 4.610 respostas de Questões Discursivas de Conhecimento Específico e 2.297 de Formação Geral.
Pelos professores corretores foram identificadas mais 52 e 18 questões em branco, respectivamente, para Conhecimento Específico e para Formação Geral.
No caso da Área de Geografia ( Bacharelado ), foram identificados como brancos “ automáticos ”, 857 respostas de Questões Discursivas de Conhecimento Específico e 475 de Formação Geral.
Pelos professores corretores foram identificadas mais nove e quatro questões em branco, respectivamente, para Conhecimento Específico e para Formação Geral.
Uma outra situação a ser mencionada é a ocorrências de atendimento a portadores de necessidades especiais.
O Anexo XI apresenta, para o exame como um todo e para a Área de Geografia ( Licenciatura/Bacharelado ), por tipo de deficiência, o protocolo usado para permitir a correção das respostas dos alunos em cada situação. 6.3.1 COMPONENTE DE FORMAÇÃO GERAL As análises dos resultados de desempenho dos estudantes de Geografia ( Licenciatura/Bacharelado ), nas duas Questões Discursivas relativas a Formação Geral, encontram -se na Tabela 6.13 e no Gráfico 6.6.
Aqui analisa -se a média das notas de cada Questão Discursiva de Formação Geral que leva em conta as correções realizadas por duas bancas, uma que analisa e pontua aspectos relativos ao conteúdo das respostas, seguindo o padrão de respostas divulgado pelo Inep ( ver Anexos IXA e IXB ), outra que avalia o desempenho linguístico expresso pela redação das respostas, seguindo um padrão de respostas específico da Área de Língua Portuguesa ( ver Anexos IXA e IXB ).
A nota de cada questão é composta de 80 \% da nota da correção de conteúdo e 20 \% da nota de desempenho linguístico.
Na Tabela 6.13, observa -se que a nota Média nesse conjunto de questões foi inferior à obtida nas objetivas pelos estudantes de Licenciatura.
Os estudantes de Licenciatura de todo o Brasil obtiveram, em Formação Geral, Média 39,6 ( Tabela 6.7a ) nas questões objetivas e 27,0 nas questões discursivas.
Pode -se notar, também, que o Desvio padrão nesse conjunto de questões foi menor do que o obtido nas objetivas : 19,7 nas questões objetivas e 17,5 nas questões discursivas.
A maior Média foi obtida na região Sudeste ( 30,6 ) e a menor, na região Sul ( 24,8 ).
A Mediana dos alunos de Licenciatura de todo o Brasil, neste componente, foi 26,5.
Na região Sudeste, a Mediana foi maior ( 30,0 ), e, nas demais regiões, foi menor que o valor para o Brasil : Norte e Nordeste ( 25,5 ), Sul ( 24,0 ) e Centro-Oeste ( 26,0 ).
A nota Máxima ( 83,5 ) foi obtida na região Sudeste, sendo 75,5, a Máxima na região Norte ; 79,0, na região Nordeste ; 82,0, na região Sul ; e 75,0, na região Centro-Oeste.
A nota Mínima ( 0,0 ) foi a mesma em todas as regiões do Brasil.
Considerando -se agora os alunos de cursos Presenciais de Licenciatura, observa -se que a nota Média nesse conjunto de questões foi inferior da obtida nas objetivas.
Estes estudantes obtiveram, em Formação Geral, Média 41,8 ( Tabela 6.7a ) nas questões objetivas e 28,6 nas questões discursivas.
Pode -se notar também que o Desvio padrão nesse conjunto de questões foi menor que o obtido nas objetivas, 19,7 nas questões objetivas e 17,7 nas questões discursivas.
A maior Média foi obtida na região Sudeste ( 32,1 ), e a menor, na região Norte ( 25,9 ).
A Mediana dos alunos de cursos Presenciais de Licenciatura, nas questões discursivas deste componente, foi 28,5.
Nas regiões Norte ( 25,5 ), Nordeste ( 27,0 ) e Centro-Oeste ( 26,0 ) a Mediana foi menor e nas demais regiões foi maior que o valor para o Brasil : Sudeste ( 32,0 ) e Sul ( 31,5 ).
A nota Máxima ( 80,0 ) foi a mesma na região Sudeste, sendo 75,5 na região Norte, 79,0 na região Nordeste, 79,5, na região Sul e 75,0, na região Centro-Oeste.
A nota Mínima ( 0,0 ) foi a mesma em todas as regiões do Brasil.
A região Norte foi a única que não apresentou cursos a Distância de Licenciatura.
A Média das notas da prova dos alunos de cursos a Distância de Licenciatura, foi 24,9, mais baixa do que os alunos de cursos Presenciais ( 28,6 ) sendo que os alunos da região Nordeste obtiveram a Média mais baixa ( 21,1 ).
O Desvio padrão das Licenciaturas a Distância foi 17,0, sendo o menor obtido na região Nordeste ( 15,1 ) e o maior na região Sudeste ( 19,0 ).
A maior Mediana foi a da região Sudeste ( 28,0 ) e a menor foi a da região Nordeste ( 20,0 ).
A maior nota Máxima foi 83,5, obtida na região Sudeste, registrada como a nota Máxima de todas as Licenciaturas.
A nota Mínima foi zero em todas as regiões.
Considerando -se agora os alunos de cursos de Bacharelado, na Tabela 6.13, observa -se que a nota Média nesse conjunto de questões foi inferior à obtida nas objetivas pelos estudantes de Bacharelado.
Os estudantes de Bacharelado de todo o Brasil obtiveram, em Formação Geral, Média 47,6 ( Tabela 6.7b ), nas questões objetivas e 32,9, nas questões discursivas.
Pode -se notar, também, que o Desvio padrão nesse conjunto de questões foi menor do que o obtido nas objetivas : 19,5, nas questões objetivas e 18,9, nas questões discursivas.
A maior Média foi obtida na região Sudeste ( 36,0 ), e a menor, na região Norte ( 25,6 ).
A Mediana dos alunos de Bacharelado de todo o Brasil, neste componente, foi 34,0.
A Mediana foi maior na região Sudeste ( 37,5 ) e, nas demais regiões, foi menor que o valor para o Brasil : Norte ( 27,0 ), Nordeste ( 33,0 ), Sul ( 32,5 ) e Centro-Oeste ( 25,0 ).
A nota Máxima ( 88,5 ) foi obtida na região Sudeste, sendo 60,0, a nota Máxima na região Norte ; 83,0, na região Nordeste ; 85,5, na região Sul ; e 74,0, na região Centro-Oeste.
A nota Mínima ( 0,0 ) foi a mesma em todas as regiões do Brasil.
Considerando -se agora os alunos de cursos Presenciais de Bacharelado, observa -se que a nota Média nesse conjunto de questões foi inferior da obtida nas objetivas.
Estes estudantes obtiveram, em Formação Geral, Média 47,9 ( Tabela 6.7b ) nas questões objetivas e 33,3 nas questões discursivas.
Pode -se notar também que o Desvio padrão nesse conjunto de questões foi menor que o obtido nas objetivas, 19,4 nas questões objetivas e 19,0 nas questões discursivas.
A maior Média foi obtida na região Sudeste ( 36,8 ), e a menor, na região Norte ( 25,6 ).
A Mediana dos alunos de cursos Presenciais de Bacharelado, nas questões discursivas deste componente, foi 35,0.
Na região Sudeste ( 38,5 ) a Mediana foi maior e, nas demais regiões, foi menor que o valor para o Brasil : regiões Norte ( 27,0 ), Nordeste ( 33,0 ), Sul ( 34,5 ) e Centro-Oeste ( 25,0 ).
A nota Máxima ( 88,5 ) foi a mesma na região Sudeste, sendo 60,0 na região Norte, 83,0 na região Nordeste, 85,5, na região Sul e 74,0, na região Centro-Oeste.
A nota Mínima ( 0,0 ) foi a mesma em todas as regiões do Brasil.
Somente nas regiões Sudeste e Sul foram oferecidos cursos a Distância de Bacharelado.
Respectivamente, para o Brasil, a região Sudeste e a região Sul as Médias foram 24,2, 20,1 e 26,3 ; o Desvio padrão foi 13,9, 14,2 e 13,3 ; a nota Mínima foi zero ; a Mediana foi 23,8, 20,0 e 25,5 ; e a nota Máxima foi 67,5, 67,5 e 59,5.
O Gráfico 6.6 representa a distribuição das notas nas Questões Discursivas do Componente de Formação Geral.
A moda desta distribuição ocorre no intervalo ( 20 ; 30 ] para os cursos de Licenciatura e no intervalo, ( 30 ; 40 ], para o Bacharelado.
Destaca -se também o intervalo [ 0 ; 10 ], que inclui, além da nota zero, a frequência de estudantes que deixaram esse tipo de questão em branco, com 15,0 \% e 19,3 \% para Bacharelados e para Licenciaturas, respectivamente. tipo de questão em branco, com 15,0 \% e 19,3 \% para Bacharelados e para Licenciaturas, respectivamente.
Na sequência, os resultados verificados para cada uma das Questões Discursivas de Formação Geral serão apresentados, estabelecendo -se relações com os conteúdos abordados em cada uma delas.
Os comentários das Bancas de docentes corretores, a respeito do observado na correção das respostas dos estudantes, suas impressões e conclusões serão apresentados junto à análise de cada questão.
Cumpre esclarecer que, tendo em vista que as questões discursivas de Formação Geral são padronizadas, ou seja, constam de todas as provas, os comentários da Banca são os mesmos para todas as carreiras acadêmicas, sendo direcionados a todos os estudantes que participaram do Enade/2021.
A seguir, serão analisados os desempenhos linguísticos dos estudantes da Área de Geografia ( Licenciatura/Bacharelado ) nas duas Questões Discursivas de Formação Geral do Enade/2021, apresentando -se os resultados obtidos e os comentários da banca de correção para cada questão. 6.3.1.1 ANÁLISE DE CONTEÚDO DA QUESTÃO DISCURSIVA 1 DO COMPONENTE DE FORMAÇÃO GERAL Os dados de Geografia ( Licenciatura/Bacharelado ), obtidos a partir das respostas à Questão 1, encontram -se na Tabela 6.14 e no Gráfico 6.7.
Nessa questão – de desempenho inferior ao da outra questão de Formação Geral –, os estudantes de Licenciatura de todo o Brasil tiveram Média 14,9.
A maior Média para a Questão 1 foi obtida na região Sudeste ( 19,3 ), e a menor, na região Sul ( 10,6 ).
Quanto à variabilidade das notas, o Desvio padrão de todo o Brasil foi 18,2.
O menor Desvio padrão foi obtido na região Sul ( 15,4 ), e o maior Desvio padrão foi obtido pela região Sudeste ( 20,2 ).
A Mediana do Brasil, para os alunos de Licenciatura, como um todo, foi 10,0, sendo igual na região Centro-Oeste, maior ( 15,0 ) nas regiões Norte, Nordeste e Sudeste, e menor na região Sul ( 0,0 ), significando que pelo menos 50 \% dos participantes da região Sul zeraram essa questão.
A nota Máxima da Questão Discursiva 1 foi, sem exceção, a mesma para todas as regiões do Brasil, 95,0.
Dentre os alunos de Bacharelado, os dados obtidos a partir das respostas à Questão 1 – de desempenho inferior ao verificado na outra questão de Formação Geral – os estudantes de Bacharelado de todo o Brasil, tiveram Média 19,4.
A maior Média para a Questão 1 foi obtida na região Sudeste ( 23,0 ), e a menor, na região Norte ( 12,0 ).
Quanto à variabilidade das notas, o Desvio padrão de todo o Brasil foi 18,9.
O menor Desvio padrão foi obtido na região Norte ( 14,0 ), e o maior Desvio padrão foi obtido pela região Sudeste ( 20,1 ).
A Mediana do Brasil, para os alunos de Bacharelado, como um todo, foi 15,0, sendo igual nas regiões Nordeste e Sul.
Na região Sudeste, a Mediana foi igual a 25,0, na região Centro-Oeste foi 5,0 e na região Norte foi zero, indicando que, na região Norte, pelo menos 50 \% dos estudantes tiraram nota zero.
A nota Máxima ( 85,0 ) da Questão Discursiva 1 foi a mesma das regiões Nordeste, Sudeste e Sul, enquanto foi 50,0 na região Norte e 60,0 na região Centro-Oeste.
A nota Mínima ( 0,0 ) foi a mesma para todas as regiões.
No Gráfico 6.7, mostra -se a distribuição das notas na Questão Discursiva 1 do Componente de Formação Geral.
Sem levar em consideração que 36,0 \% dos alunos de Licenciatura receberam nota zero, a moda dessa distribuição seria o intervalo ( 10 ; 20 ], com 19,8 \% do total de participantes.
Dentre os alunos de Bacharelado, sem levar em consideração que 23,3 \% receberam nota zero, a moda dessa distribuição seria o intervalo ( 20 ; 30 ], com 22,8 \% do total de participantes. 6.3.1.2 COMENTÁRIOS SOBRE A CORREÇÃO DE CONTEÚDO DAS RESPOSTAS À QUESTÃO DISCURSIVA 1 O enunciado da Questão1 continha dois textos e um comando.
O primeiro texto trazia uma abordagem sobre as tensões entre cultura, arte e censura sob o recorte do inciso IX do Art. 5º da Constituição Federal de 1988, que diz ser livre a expressão da atividade intelectual, artística, científica e de comunicação, independentemente de censura ou licença.
O segundo texto apresentava o inciso IX do parágrafo 5º da Constituição Federal, em que é garantida a livre expressão da atividade artística, científica e de comunicação, independentemente de censura ou licença.
O comando solicitava que o respondente discorresse a respeito de dois aspectos.
Primeiramente, o estudante deveria ser capaz de dissertar sobre as relações entre os conceitos de arte, cultura e censura, à luz da ideia de liberdade artística garantida pela Constituição Federal de 1988.
O estudante também deveria ser capaz de apresentar duas ações educativas que pudessem contribuir “ para minimizar essas tensões ” ( tensões entre arte, cultura e censura ) e garantir a liberdade artística prevista pela lei.
A questão, portanto, estava de acordo com os temas previstos nas Diretrizes Curriculares para a Formação Geral : ética, democracia e cidadania ; cultura, arte e comunicação ; educação e desenvolvimento humano e social.
Além disso, o comando permitia avaliar o desenvolvimento das seguintes competências previstas nas Diretrizes Curriculares para a Formação Geral : planejar e elaborar projetos de ação e intervenção a partir da análise de necessidades em contextos diversos ; ler, produzir e interpretar textos com clareza e coerência ; identificar, compreender e analisar situações-problema utilizando pensamento lógico e sistêmico, estabelecendo relações de \textbf{causalidade}.
O enunciado da questão era claro, empregava linguagem adequada e continha informações suficientes para que os egressos construíssem seus argumentos.
Entretanto, apesar dessas características favoráveis, a temática era sofisticada, exigia reflexão e estabelecimento de relações a partir de textos concisos e densos.
Além disso, os conceitos de arte e cultura expressos pelo Texto I excluíam inúmeras outras definições reconhecidas em diferentes áreas das ciências humanas, em alguns casos divergentes daquelas apresentadas.
Os conhecimentos e o nível de profundidade exigidos pelo padrão de resposta eram adequados para egressos de cursos superiores.
Deve -se destacar que o padrão admitia um espectro amplo para a construção da argumentação das respostas, viabilizando a aceitação de outras interpretações dos conceitos de arte e cultura, incluindo possíveis contestações em relação às ideias apresentados no Texto I.
A questão pode ser considerada difícil, já que a primeira parte exigia uma resposta que articulasse quatro conceitos complexos – arte, cultura, censura e liberdade artística –, e a segunda parte exigia o conhecimento do que seja uma ação educativa, além da proposição de ações que objetivamente contribuíssem para minimizar as tensões entre arte, cultura e censura, além de garantir a liberdade artística.
Essas \textbf{observações}, aliadas ao fato do espaço e do tempo limitados para elaboração de uma resposta tão extensa, elevaram consideravelmente o nível de dificuldade da questão.
A correção confirmou a dificuldade dos estudantes para abordar a temática, evidenciou grande diversidade na compreensão dos conceitos envolvidos e falta de clareza do que seja uma ação educativa.
Nas respostas à primeira parte do comando, os estudantes raramente abordaram os quatro conceitos que deveriam estar presentes na resposta.
A maioria se limitou a relacionar arte e cultura com comentários que repetiam trechos dos enunciados, como arte e cultura"andam de mãos dadas".
Assim, apesar de o Texto I sugerir o entendimento de que a cultura é uma esfera reguladora do comportamento e de que a arte é o elemento de desequilíbrio, de contestação, a quase totalidade das respostas não fez uso dessa interpretação, tratando arte e cultura quase como sinônimos.
Apesar de raras, houve respostas nas quais os estudantes defenderam a existência de múltiplas culturas, muitas delas não hegemônicas, tão suscetíveis aos dispositivos de controle quanto a arte.
Houve também respostas que colocavam em questão os critérios, muitas vezes elitistas e eurocêntricos, utilizados para produzir hierarquização entre as artes, chegando a inviabilizar o reconhecimento de outras tantas manifestações artísticas.
Foi pequeno o grupo que, além de relacionar arte e cultura, foi capaz de articulá -las com o conceito de censura trazendo -a, quase sempre, como herança do período militar.
Entretanto, nesse grupo, outras abordagens de censura foram mapeadas no decorrer das correções : cancelamento em redes sociais ; contrainformação ou manipulação dos meios de comunicação social ; além de outros recursos censórios mais indiretos, como os processos judiciais, a suspensão de financiamentos públicos a projetos culturais e a autocensura.
Em relação à segunda parte do comando, era esperado que as/os estudantes egressas/os citassem duas ações educativas sob o recorte de que minimizassem tensões observadas entre arte, cultura e censura e garantissem a liberdade artística prevista pela lei.
Na maioria das respostas, não se verificou essa relação entre a proposição e seu objetivo/resultado ; muitas ações foram apenas citadas sem explicações.
Esse fato reduziu bastante a quantidade de respostas com pontuação integral no referido item.
Importante ressaltar que o padrão de resposta indicou exemplos de ações educativas no campo da educação não formal.
Entretanto, por não ser algo explícito no comando, foram também aceitas ações referentes ao campo da educação formal.
O mais recorrente nas respostas foi a indicação de ações educativas em sala de aula e visitas a museus e a exposições, mediadas pelo próprio professor da turma.
Entretanto, houve respostas em que os estudantes ampliaram tanto o espaço quanto o grupo considerado nas ações educacionais, propondo campanhas para disseminar conhecimento sobre arte e/ou cultura para plateia maior do que a de estudantes.
De modo geral, a fluidez do texto das respostas e a consistência das ideias articuladas ficou aquém do padrão desejado para estudantes egressos do ensino superior.
O maior volume de respostas mostrou que os estudantes valorizaram mais a apresentação de duas ações educativas, muitas vezes de forma incompleta, apenas citando uma ação proposta e deixando de explicar o desdobramento solicitado no comando.
Além disso, houve respostas em que sugeriram, de forma isolada, ações do tipo apresentações, shows, exposições, sem qualquer detalhamento. shows, exposições, sem qualquer detalhamento.
Entretanto, uma ação educativa com as finalidades solicitadas precisaria ser minimamente descrita para que fosse possível identificar que a ação pressupunha a presença de um mediador ( professor / arte educador / museólogo ) e de receptores participativos.
A proposição de uma ação educativa também foi confundida pelos estudantes com a proposta de implantação de equipamentos de infraestruturas ou projeto governamental, tais como : escolas, centros educacionais, incentivo ao ensino das artes em áreas mais necessitadas e introdução de programas de ensino de artes nos currículos escolares.
O contexto da arte e/ou da cultura e seus confrontos com a censura foram considerados, muitas vezes, sem o estabelecimento de relações, mas restringindo -se a aspectos descritivos de como o respondente compreendia cada um desses elementos.
Há de se considerar, adicionalmente, que em uma parte efetiva das respostas os estudantes defendiam, indevidamente, a censura como um mecanismo legitimamente regulador dos chamados"excessos da arte", revelando confusão entre os limites legais estabelecidos em lei ( que tipificam discursos de ódio, pedofilia, apologia às drogas, etc. ) e os julgamentos pessoais, coletivos ou institucionais que \textbf{desencadeiam} atitudes censoras.
A abrangência proporcionada pelo tema e pelo comando refletiu -se nas respostas, principalmente nos exemplos das ações educativas apresentados pelos estudantes.
Já na reflexão sobre arte, cultura e censura não houve tanta diversidade – a discussão quase sempre permanecia no senso comum, com menção à influência da época militar ou associada à discussão atual sobre limitações à liberdade de expressão.
Houve grande incidência de respostas em que o elemento censura foi confundido com a noção de classificação indicativa.
O desempenho ficou aquém do esperado.
Na maioria das respostas, os alunos ficaram com notas baixas, e foi grande a incidência de zeros.
Cabe ressaltar também a grande quantidade de respostas em branco e as que receberam nota zero devido a anulações por defesa de atos de censura, na contramão do que estabelecem a Constituição Federal e os Direitos Humanos.
A fuga ao que foi solicitado pelo comando da questão foi marcante, talvez evidenciando dificuldade de interpretação ou distanciamento em relação ao tema.
Outra lacuna muito observada foi a dificuldade de se expressar em textos fluentes e articulados.
A fragmentação e a falta de profundidade da argumentação foram uma constante.
Além disso, chamou a atenção a predominância do lugar comum na construção das respostas.
Tratando -se de um universo de formandos de nível superior, a superficialidade da reflexão em torno da arte e da censura e a dificuldade para incluir a Constituição Federal na discussão foram surpreendentes. 6.3.1.3 ANÁLISE DE CONTEÚDO DA QUESTÃO DISCURSIVA 2 DO COMPONENTE DE FORMAÇÃO GERAL A Tabela 6.15, mostra que o desempenho médio dos estudantes de Licenciatura, na Questão Discursiva 2 ( Média 29,7 ), foi próxima ao obtido na Questão Discursiva 1 ( Média 14,9 ).
A região Norte foi aquela cuja Média foi menor ( 27,5 ), e a maior Média foi na região Sudeste ( 33,2 ).
Quanto à variabilidade das notas, o Desvio padrão de todo o Brasil, entre os cursos de Licenciatura, foi 27,0, superior ao obtido na Questão Discursiva 1 ( 18,2 ).
O maior desvio nessa questão foi obtido na região Sudeste ( 28,2 ), enquanto o menor foi obtido na região Centro-Oeste ( 25,2 ).
A Mediana de todo o Brasil, entre os cursos de Licenciatura, foi 25,0, a mesma de quatro regiões.
A exceção foi a região Sudeste, que obteve Mediana 30,0.
A nota Máxima ( 100,0 ) foi a mesma nas regiões Nordeste, Sudeste e Sul, e nas demais regiões foi 95,0.
A nota Mínima ( 0,0 ) foi a mesma em todas as regiões.
Dentre os alunos de Bacharelado, o desempenho médio dos estudantes, na Questão Discursiva 2 ( 38,9 ), foi superior ao obtido na Questão Discursiva 1 ( 19,4 ).
A região Centro-Oeste foi aquela cuja Média, também nessa questão, foi menor ( 31,4 ), mas a de maior Média nessa questão foi a região Sudeste ( 42,0 ).
Quanto à variabilidade das notas, o Desvio padrão de todo o Brasil, entre os cursos de Bacharelado, foi 29,7, superior ao obtido na Questão Discursiva 1 ( 18,9 ).
O maior desvio nessa questão foi obtido na região Sudeste ( 30,0 ), enquanto o menor foi obtido na região Centro-Oeste ( 29,1 ).
A Mediana de todo o Brasil, entre os cursos de Bacharelado, foi 40,0, a mesma das regiões Nordeste e Sul, enquanto foi 30,0 nas regiões Norte e Centro-Oeste, e 45,0 na região Sudeste.
As notas Máxima ( 100,0 ) e Mínima ( 0,0 ) foram as mesmas em todas as regiões, sem exceção.
No Gráfico 6.8, mostra -se a distribuição das notas na Questão Discursiva 2 do Componente de Formação Geral.
Observa -se que 11,6 \% dos estudantes de Licenciatura deixaram a questão em branco.
Chega a 28,2 \%, se agregado à frequência daqueles que, tendo respondido, receberam nota zero.
O intervalo modal foi o ( 20 ; 30 ], com 19,2 \% do total de participantes.
Nota -se, ainda, que as notas ficaram mais dispersas em comparação às da Questão Discursiva de número 1, o que pode ser constatado, também, pela comparação do Desvio padrão das notas da Questão Discursiva 2 ( 27,0 ) e o das notas da Questão Discursiva 1 ( 18,2 ).
Dentre os alunos de Bacharelado, observa -se que 11,6 \% dos estudantes deixaram a resposta à questão em branco.
Chega a 21,5 \%, se agregado à frequência daqueles que, tendo respondido, receberam nota zero, a classe modal.
O intervalo modal foi o ( 20 ; 30 ], com 17,1 \% No Gráfico 6.8, mostra -se a distribuição das notas na Questão Discursiva 2 do Componente de Formação Geral.
Observa -se que 11,6 \% dos estudantes de Licenciatura deixaram a questão em branco.
Chega a 28,2 \%, se agregado à frequência daqueles que, tendo respondido, receberam nota zero.
O intervalo modal foi o ( 20 ; 30 ], com 19,2 \% do total de participantes.
Nota -se, ainda, que as notas ficaram mais dispersas em comparação às da Questão Discursiva de número 1, o que pode ser constatado, também, pela comparação do Desvio padrão das notas da Questão Discursiva 2 ( 27,0 ) e o das notas da Questão Discursiva 1 ( 18,2 ).
Dentre os alunos de Bacharelado, observa -se que 11,6 \% dos estudantes deixaram a resposta à questão em branco.
Chega a 21,5 \%, se agregado à frequência daqueles que, tendo respondido, receberam nota zero, a classe modal.
O intervalo modal foi o ( 20 ; 30 ], com 17,1 \% dos participantes.
Nota -se, ainda, que as notas ficaram levemente mais dispersas em comparação às da Questão Discursiva de número 1, o que pode ser constatado, também, pela comparação do Desvio padrão das notas da Questão Discursiva 2 ( 29,7 ) e o das notas da Questão Discursiva 1 ( 18,9 ). 6.3.1.4 COMENTÁRIOS SOBRE A CORREÇÃO DE CONTEÚDO DAS RESPOSTAS À QUESTÃO DISCURSIVA 2 O enunciado da segunda Questão Discursiva de Formação Geral era composto por dois textos motivadores e um comando com dois itens.
O primeiro texto apresentava uma definição de cidade inteligente, pelo uso da tecnologia nos diversos setores da administração e seu compromisso com o meio ambiente.
O segundo texto ampliava a definição anterior explorando aspectos da gestão pública.
O enunciado era claro, empregava linguagem adequada e continha informações suficientes para que os egressos construíssem seus argumentos.
O comando solicitava, primeiramente, que o estudante explicasse de que modo as cidades inteligentes podem contribuir para a melhoria das questões relacionadas ao desenvolvimento sustentável.
Em atendimento ao segundo item do comando, o estudante deveria apresentar uma proposta de intervenção urbana que gerasse impacto social e contribuísse para a melhoria da vida em comunidade.
Os dois itens do comando exigiam articulação de ideias e exposições com desdobramentos do porquê da resposta, o que tornou a questão complexa.
A questão abordava, portanto, temas previstos nas Diretrizes Curriculares para a Formação Geral : ciência, tecnologia e inovação ; biodiversidade, sustentabilidade e intervenção humana e cidades, habitação e qualidade de vida.
Além disso, avaliava o desenvolvimento de competências também previstas nas Diretrizes Curriculares para a Formação Geral : planejar e elaborar projetos de ação e intervenção a partir da análise de necessidades em contextos diversos ; ler, produzir e interpretar textos com clareza e coerência ; identificar, compreender e analisar situações-problema utilizando pensamento lógico e sistêmico, estabelecendo relações de \textbf{causalidade}.
Os conhecimentos e o nível de profundidade exigidos foram considerados adequados para egressos de cursos superiores pelos componentes da banca de correção.
Os textos tratavam de temas atuais e de amplo debate na mídia, tais como : cidade inteligente, meio ambiente, desenvolvimento sustentável, intervenção urbana e impacto social.
No entanto, o nível de dificuldade foi considerado elevado, já que tais conceitos admitem interpretações diversas nesse universo particular de egressos, o que criou, em consequência, um espectro diversificado de respostas e tipos de argumentação.
Além disso, era necessário conhecer o conceito de intervenção urbana para propor uma ação deste tipo que gerasse impacto social.
O padrão de resposta estava redigido de forma clara e, para o item ‘ a ’ do comando, apresentava um exemplo simples e informativo de como a tecnologia pode ser usada para a redução do consumo de energia e emissão de CO2.
Para o item ‘ b ’, o padrão apresentava exemplos de propostas de intervenção de forma mais diversificada do que a resposta padrão do item ‘ a ’, levantando várias opções corretas.
A abordagem adotada pelos estudantes foi quase sempre a da exemplificação, o mesmo recurso adotado no padrão de resposta.
Assim, pode -se dizer que as respostas dos estudantes se mantiveram, quase sempre, em conformidade com o contexto do padrão, mas com uma variedade muito grande de temas e ambientes relacionados ao desenvolvimento sustentável.
Dentre as ações vinculadas às cidades inteligentes, as ocorrências mais significativas foram as relacionadas à poluição, ao desmatamento, aos transportes coletivos, à segurança e ao desenvolvimento de tecnologia de modo a contribuir para a melhoria das dificuldades enfrentadas atualmente nas cidades.
Em muitos casos, o estudante articulou a sua resposta com a realidade vivenciada em determinadas cidades, comparando, inclusive, cidades que são consideradas inteligentes com realidades locais.
Também foram encontradas respostas em que os estudantes apresentavam propostas e soluções para realidades específicas das cidades onde o respondente reside, contextualizando o uso da tecnologia com ações fundamentais que faltam em algumas cidades brasileiras, como saneamento básico, água não poluída para uso doméstico, destinação do lixo.
As respostas incompletas ocorreram, muitas vezes, em função de o aluno apenas apresentar ações ou características das cidades inteligentes sem articular com o desenvolvimento sustentável ou com a intervenção urbana.
Em muitas respostas, os alunos apenas repetiram ideias do enunciado ; em outras, demonstraram a não compreensão sobre o que estava sendo solicitado pelo comando, deixando de argumentar sobre uma proposta ou ação mencionada.
Nas respostas, os estudantes se concentraram no uso da tecnologia para monitorar o desmatamento e a violência ; em hortas comunitárias para melhorar a qualidade da alimentação das pessoas ; em investimento em locais públicos para a realização de exercícios ; em maior controle do uso de combustíveis não fósseis para evitar a poluição.
As propostas de intervenção fixaram -se bastante em aplicativos, solução sugerida pelo ambiente atual de desenvolvimento e das facilidades proporcionadas pela Internet.
Marcante foi a diversidade de funções encontradas nas propostas de sensores e aplicativos, desde marcador de horário de passagem de ônibus até controlador de desperdício de algum bem, como água, por exemplo.
O resultado ficou abaixo do esperado.
A quantidade expressiva de zeros talvez tenha sido consequência de desinteresse mais do que por dificuldades para interpretar e responder ao comando.
Essas ficaram mais evidentes nas propostas de intervenção urbana, muitas vezes incompletas por sequer explicitarem o benefício delas advindo.
As maiores fragilidades reveladas pelas respostas dos estudantes, apontadas pelos corretores, foram a interpretação do contexto e do comando da questão, além de uma grande dificuldade para se expressarem de forma fluente e concisa.
As respostas com geração de ideias próprias ocorreram em escala bem menor do que aquelas em que os estudantes usaram do recurso da paráfrase pobre sobre o texto do enunciado, revelando dificuldade para criar e articular ideias a respeito dos temas em foco.
A linguagem utilizada nas respostas ficou, muitas vezes, aquém do padrão desejado para estudantes egressos do ensino superior.
A clareza das ideias apresentadas e a articulação do conteúdo com o comando da questão raramente foram encontradas nas respostas no nível esperado.
A dificuldade mais evidente dos estudantes foi a de estabelecer relações e encadear ideias para completar um raciocínio iniciado ; foram muitas as respostas fragmentadas. 6.3.1.5 CONSIDERAÇÕES FINAIS SOBRE A CORREÇÃO DE CONTEÚDO DAS QUESTÕES DISCURSIVAS DO COMPONENTE DE FORMAÇÃO GERAL As Questões Discursivas da Área de Formação Geral abordaram temas atuais e proporcionaram a avaliação de competências importantes dentre as previstas pelas suas Diretrizes Curriculares.
A análise das respostas revelou, principalmente, dificuldade dos estudantes na interpretação dos enunciados – talvez pela linguagem dos textos motivadores ser mais complexa do que a do cotidiano – e na compreensão dos comandos das questões – que pode ser fruto de desatenção ou de desconhecimento mais profundo do contexto.
Esses fatos parecem ter sido os principais responsáveis pelo resultado aquém do esperado em termos de avaliação quantitativa.
Mesmo os estudantes que responderam em conformidade com o padrão demonstraram um universo de argumentação muito limitado, pobre em vocabulário, com o recurso frequente de procurar nos textos introdutórios os elementos para a construção de suas frases.
Isso criou uma peculiaridade interessante : grupos de respostas muito parecidas.
Encadear e relacionar argumentos também foi uma deficiência percebida.
As razões eram apresentadas de forma confusa, sem fluência e sem constituir unidade.
O conteúdo das respostas aponta uma preocupação com o viés social dos problemas abordados ; são bons exemplos as propostas bastante diversificadas de melhoria de vida para a parcela da população menos favorecida.
Por outro lado, traz indícios de deficiências consideráveis na separação conceitual de cultura e arte, na compreensão mais ampla do conceito de censura e na capacidade de separar ações educativas de culturais ; em algumas respostas, também se evidenciaram conceitos equivocados sobre tecnologia, com exemplos de concretização muito difícil física ou financeiramente.
Houve alguns casos em que os conceitos de sustentabilidade e de meio ambiente se confundiam em lugar de se complementarem.
O recorte do grupo de formandos avaliado, em resumo, é de dificuldade de interpretação e produção de textos, vocabulário restrito ao praticado cotidianamente, assunção de fraseologia de senso comum em lugar do conteúdo apreendido formalmente e falta de profundidade em matérias importantes para a formação do cidadão como cultura, arte, educação, liberdade de expressão, tecnologia, cidade inteligente e desenvolvimento sustentável. 6.3.1.6 ANÁLISE DE LÍNGUA PORTUGUESA DAS QUESTÕES DISCURSIVAS DO COMPONENTE DE FORMAÇÃO GERAL Os dados obtidos a partir das respostas à Questão Discursiva do Componente de Formação Geral de Geografia desagregados por Grande Região, segundo Habilitação e Modalidade de Ensino, no que tange à Língua Portuguesa, encontram -se na Tabela 6.16 e no Gráfico 6.9.
Nos cursos de Licenciatura em Geografia, a Média desse aspecto foi 46,2, inferior à obtida pelos alunos de Bacharelado ( 48,2 ).
Quando se considera a Modalidade de Ensino para cada Habilitação, os alunos de cursos Presenciais de Licenciatura obtiveram Média 46,7.
A Média com respeito à Língua Portuguesa foi menor para os alunos de cursos a Distância ( 45,6 ).
Para os cursos de Bacharelado, os valores foram : 48,1 para os cursos Presenciais, e 49,2 para os cursos a Distância.
A Mediana das notas de Língua Portuguesa tanto para os cursos de Licenciatura quanto de Bacharelado foi 52,5 para os cursos Presenciais e 50,0 para os cursos a Distância.
A nota Máxima não chegou a 100,0 para nenhuma combinação de Habilitação e Modalidade.
No Gráfico 6.9, mostra -se a distribuição das notas de Língua Portuguesa das Questões Discursivas do Componente de Formação Geral.
Observa -se que a maior frequência 28,3 \%, corresponde à dos estudantes de Licenciatura que obtiveram nota no intervalo ( 50 ; 60 ].
Destacam -se, também, os estudantes de Licenciatura que deixaram ambas as questões discursivas em branco, representando, 8,5 \% do total.
Dentre os alunos de Bacharelado, mostra -se a distribuição das notas de Língua Portuguesa do Componente de Formação Geral da questão.
Observa -se que a maior frequência ( 26,7 \% ) corresponde à dos estudantes de Bacharelado que obtiveram nota no intervalo ( 50 ; 60 ].
Destacam -se, também, os estudantes de Bacharelado que deixaram ambas as questões discursivas em branco, representando 8,3 \% do total. 6.3.1.7 COMENTÁRIOS SOBRE A CORREÇÃO DE LÍNGUA PORTUGUESA DAS RESPOSTAS ÀS QUESTÕES DISCURSIVAS A avaliação de desempenho linguístico De acordo com a portaria nº 386 do MEC/INEP, de 23 de agosto de 2021, o Enade deve avaliar se o estudante desenvolveu, no processo de formação, competências para : compreender as linguagens e respectivas variações ( item V ) ; ler, produzir e interpretar textos com clareza e coerência ( item VI ) ; analisar e interpretar representações verbais, não verbais de fenômenos diversos ( item VII ) ; identificar diferentes representações de um mesmo significado ( item VIII ) ; formular e articular argumentos e contra-argumentos consistentes em situações sociocomunicativas ( item IX ).
Essas competências são essenciais à formação de nível superior, porque dizem respeito ao domínio da linguagem, tanto no processo de leitura quanto na articulação e na organização de ideias para a produção textual.
Dessa forma, as duas questões discursivas de formação geral oferecem oportunidade para avaliar se o participante consegue desenvolver um tema, recorrendo a argumentos que explicitem sua opinião e elaborando propostas de intervenção social adequadas a um exercício consciente da cidadania, conforme determina a referida portaria : “ As questões discursivas do componente de Formação Geral avaliam aspectos como clareza, coerência, coesão, estratégias argumentativas, vocabulário e gramática adequados à norma-padrão da Língua Portuguesa. ” ( Art.2 §2 da referida portaria nº 386 ) Para avaliar esses aspectos, portanto, é necessário identificar, nas respostas discursivas de Formação Geral, as diferentes habilidades necessárias ao bom desempenho linguístico na situação de formalidade exigida : a organização das ideias, tendo em vista a construção de um texto coerente e estruturado ; o domínio das convenções ortográficas e das convenções morfossintáticas que constituem o arcabouço da gramática da norma-padrão da Língua Portuguesa ; e o domínio de um vocabulário amplo e diversificado, que permita uma comunicação clara e eficiente.
A Questão Discursiva 1 desenvolveu o tema da relação entre arte, cultura e censura, tomando como base dois fragmentos de textos.
O enunciado estava claro ao encaminhar a solicitação de resposta, evidenciando a exigência de elaboração de um texto dissertativo, por meio do verbo"discorrer", e de apresentação de duas ações educativas, articuladas com o ponto de vista desenvolvido na parte argumentativa.
A contraposição dos dois textos favoreceu a formulação e a articulação de argumentos e contra-argumentos consistentes, em defesa do ponto de vista a ser defendido pelo participante.
O enunciado permitia, inclusive, a ampliação da abordagem com base na análise de situações concretas que evidenciassem essa articulação.
A Questão Discursiva 2 desenvolveu o conceito de"cidade inteligente", tomando como base dois fragmentos de estudos que descreviam as suas características.
Com base nesses dois fragmentos, o comando solicitou do participante :"explique a contribuição das cidades inteligentes para o desenvolvimento sustentável"e apresente uma proposta de intervenção urbana que contribua para a melhoria da vida em comunidade.
O comando “ explicar ” demandava a elaboração de um texto de base expositiva, sem desenvolvimento de argumentação favorável ou contrária à configuração de uma cidade inteligente e de seus benefícios para a sociedade.
Em relação à segunda parte, o comando exigia a elaboração de uma proposta de intervenção urbana, que deveria conter informações sobre os agentes da intervenção, os recursos aplicados e o detalhamento dessa ação.
A segunda questão, portanto, se diferenciava da primeira, por solicitarem estruturas textuais distintas : a Questão 1 se consolidava com base em um modo de organização argumentativo e a Questão 2 se consolidava com base em um modo de organização expositivo.
Esses dois estilos textuais ofereceram oportunidade de avaliar se os participantes estavam habilitados a desenvolverem textos com esses dois distintos objetivos comunicativos, o que seria esperado ao final do curso de graduação e necessário à sua vida profissional.
A análise do desempenho linguístico esperado de um estudante que está concluindo seu curso superior apoiou -se no princípio de que ele deveria utilizar conhecimentos sobre os temas e estruturar seus textos de acordo com as características do registro formal adequado à situação comunicativa – avaliação de conhecimentos ao final de um curso superior.
Essa configuração determina exigências quanto à adequação da seleção vocabular, ao desenvolvimento do conteúdo, à estruturação sintática dos períodos, à organização lógica das ideias, à utilização de procedimentos de encadeamento textual e de referenciação, à obediência às exigências morfossintáticas próprias da modalidade escrita da norma-padrão, ao respeito às regras ortográficas de acentuação gráfica e de grafia das palavras.
O padrão de resposta utilizado na avaliação A partir desses preceitos, foram considerados, no padrão de resposta, as seguintes competências : ( 1 ) estruturação textual ; ( 2 ) respeito às convenções ortográficas da norma-padrão da Língua Portuguesa ; ( 3 ) domínio dos diferentes aspectos morfossintáticos próprios da modalidade escrita formal da norma-padrão da Língua Portuguesa ; e ( 4 ) seleção vocabular adequada à modalidade escrita formal da Língua Portuguesa, exigida pela situação comunicativa.
A seguir, os aspectos avaliados por cada uma destas competências, são apresentados. ( 1 ) Estruturação textual condizente com o gênero textual solicitado e o modo de organização textual adequado ao gênero ( argumentativo ou expositivo ).
Esta competência envolve : - estruturação sintática condizente com o padrão da modalidade escrita formal da Língua Portuguesa de modo a garantir a clareza e a coerência necessárias ; - distribuição do conteúdo do texto em parágrafos, de modo a garantir a sua organização temática ; - distribuição das ideias em períodos que mantenham uma estruturação sintática condizente com a modalidade formal da língua ( evitando -se frases fragmentadas, sequência justaposta de ideias sem encaixamentos sintáticos e presença de oração subordinada sem oração principal que reflitam hábitos da oralidade e comprometam a estrutura lógico-gramatical do texto ) ; - utilização de operadores discursivos que contribuam para a progressão temática do texto, estabelecendo relações lógicas entre as ideias apresentadas, tanto do ponto de vista intrafrasal, como interfrasal : preposições, conjunções, pronomes relativos, advérbios e locuções adverbiais de tempo e lugar ; - utilização de procedimentos de referenciação lexical e pronominal que permitam a retomada de referentes textuais, como pronomes, advérbios e artigos ; além de elementos lexicais, como sinônimos, hiperônimos, hipônimos, expressões metafóricas e expressões resumitivas ) ; - emprego adequado dos pronomes relativos “ cujo(a) ” e “ onde ”, responsáveis por encaixamentos sintáticos no período composto ; - utilização adequada de sinais de pontuação que contribuam para a organização lógica da frase e do texto ; - inteligibilidade relacionada ao atendimento das exigências de estruturação textual.
A expectativa, portanto, era de que o participante recorresse a procedimentos linguístico-discursivos para organizar seu texto, permitindo o encadeamento lógico entre suas partes de forma a garantir a progressão e a coerência textuais, com estruturação que não reproduzisse hábitos da oralidade e sim atendesse à norma-padrão da Língua Portuguesa. ( 2 ) Respeito às convenções ortográficas da norma-padrão da Língua Portuguesa.
Esta competência envolve o respeito às convenções estabelecidas pela legislação em vigor e consubstanciadas no Vocabulário Ortográfico da Língua Portuguesa, editado pela Academia Brasileira de Letras e avaliam -se : - domínio das regras de acentuação gráfica ; - domínio da grafia padrão das palavras ( com ausência de abreviaturas próprias da linguagem da internet ).
Assim, esperava -se que o participante : grafasse corretamente as palavras ; respeitasse as regras de acentuação gráfica ; empregasse maiúsculas em início de frase, em nomes próprios de pessoas, lugares ou instituições ; grafasse as siglas com letras maiúsculas ou apenas com a primeira letra maiúscula quando formarem uma palavra ; evitasse abreviações como p/, vc, tb, pq, tá, né, usadas muitas vezes em escrita informal e na internet ; obedecesse às regras de separação de sílabas no final da linha ; atendesse aos padrões de segmentação vocabular da norma-padrão. ( 3 ) Domínio dos diferentes aspectos morfossintáticos próprios da modalidade escrita formal da norma-padrão da Língua Portuguesa.
Esta competência envolve o domínio das convenções próprias da norma padrão da Língua Portuguesa relativas aos seguintes aspectos : concordância nominal ; concordância verbal ; regência nominal ; regência verbal ; flexão nominal ; flexão verbal ; correlação entre os tempos verbais ; colocação pronominal. vista intrafrasal, como interfrasal : preposições, conjunções, pronomes relativos, advérbios e locuções adverbiais de tempo e lugar ; - utilização de procedimentos de referenciação lexical e pronominal que permitam a retomada de referentes textuais, como pronomes, advérbios e artigos ; além de elementos lexicais, como sinônimos, hiperônimos, hipônimos, expressões metafóricas e expressões resumitivas ) ; - emprego adequado dos pronomes relativos “ cujo(a) ” e “ onde ”, responsáveis por encaixamentos sintáticos no período composto ; - utilização adequada de sinais de pontuação que contribuam para a organização lógica da frase e do texto ; - inteligibilidade relacionada ao atendimento das exigências de estruturação textual.
A expectativa, portanto, era de que o participante recorresse a procedimentos linguístico-discursivos para organizar seu texto, permitindo o encadeamento lógico entre suas partes de forma a garantir a progressão e a coerência textuais, com estruturação que não reproduzisse hábitos da oralidade e sim atendesse à norma-padrão da Língua Portuguesa. ( 2 ) Respeito às convenções ortográficas da norma-padrão da Língua Portuguesa.
Esta competência envolve o respeito às convenções estabelecidas pela legislação em vigor e consubstanciadas no Vocabulário Ortográfico da Língua Portuguesa, editado pela Academia Brasileira de Letras e avaliam -se : - domínio das regras de acentuação gráfica ; - domínio da grafia padrão das palavras ( com ausência de abreviaturas próprias da linguagem da internet ).
Assim, esperava -se que o participante : grafasse corretamente as palavras ; respeitasse as regras de acentuação gráfica ; empregasse maiúsculas em início de frase, em nomes próprios de pessoas, lugares ou instituições ; grafasse as siglas com letras maiúsculas ou apenas com a primeira letra maiúscula quando formarem uma palavra ; evitasse abreviações como p/, vc, tb, pq, tá, né, usadas muitas vezes em escrita informal e na internet ; obedecesse às regras de separação de sílabas no final da linha ; atendesse aos padrões de segmentação vocabular da norma-padrão. ( 3 ) Domínio dos diferentes aspectos morfossintáticos próprios da modalidade escrita formal da norma-padrão da Língua Portuguesa.
Esta competência envolve o domínio das convenções próprias da norma padrão da Língua Portuguesa relativas aos seguintes aspectos : concordância nominal ; concordância verbal ; regência nominal ; regência verbal ; flexão nominal ; flexão verbal ; correlação entre os tempos verbais ; colocação pronominal.
Assim, esperava -se que o participante : - flexionasse o verbo para estabelecer concordância de número com o sujeito da frase ; - flexionasse o artigo, o adjetivo e o pronome para concordar em número e em gênero com o substantivo a que cada um se refere ; - observasse a regência nominal e a verbal, utilizando a preposição adequada depois de um substantivo, um verbo ou um adjetivo ; - empregasse adequadamente o acento grave indicador da crase que expresse a contração entre uma preposição e um artigo, de acordo com as exigências de regência nominal ou verbal ; - obedecesse às regras de colocação pronominal ( próclise e ênclise ), distintas dos hábitos da oralidade ou da escrita informal ; - flexionasse adequadamente verbos, substantivos, adjetivos e pronomes, no que diz respeito à expressão das categorias gramaticais ; - flexionasse os verbos para expressar a correlação de modo e de tempo nas estruturas subordinadas.
Com base nesses critérios, alguns desvios ortográficos foram considerados como desvios de caráter morfossintático porque interferem na classificação dos vocábulos de base morfológica e funcional : - eliminação da marca de infinitivo ( -r- ) e substituição por acento agudo ou ausência total de marca do infinitivo, causando ambiguidade entre distintas classes de palavras ; - confusão entre as terminações “ ão ” e “ am ”, utilizadas para marcar diferentes tempos verbais ; - emprego indevido da forma verbal de"haver"\&lt ; “ há”\&gt ; no lugar da preposição ou do artigo \&lt ; “ a”\&gt ;, e vice-versa ; - emprego de hífen para separar pronome átono – tanto uso indevido quanto omissão ( ex : “ esperasse ”, em lugar de “ espera -se ” ; “ falar -mos ”, no lugar de “ falarmos ” ) ; - emprego de “ esta ” no lugar de “ está ”, provocando ambiguidade devido à confluência gráfica entre o pronome demonstrativo e a forma do verbo"estar", pela ausência da acentuação gráfica ; - emprego da conjunção aditiva"e"no lugar da forma verbal “ é ”, provocando ambiguidade devido à confluência gráfica por ausência da acentuação gráfica ; - emprego do advérbio “ mais ” no lugar da conjunção adversativa “ mas ”, reproduzindo confluência fônica da oralidade ; - registro da terceira pessoa do plural dos verbos “ ter ” e “ vir ” sem o acento circunflexo, causando desvio de concordância verbal ; - emprego inadequado do relativo “ onde ” deslocado do seu sentido locativo, reproduzindo tendência da oralidade ; - emprego do pronome relativo “ que ” sem a necessária preposição ( desvio de regência ) ; - ausência do acento grave obrigatório para marcar a crase. ( 4 ) Seleção vocabular adequada à modalidade escrita formal da Língua Portuguesa, exigida pela situação comunicativa.
Esta competência envolve : - precisão na seleção/utilização do vocabulário relacionado à temática solicitada pela questão ; - diversificação e adequação vocabular ; - coerência na combinação vocabular ; - ausência de marcas da oralidade, como termos de sentido muito genérico ( “ coisa ”, “ negócio ”, “ você ” ) e termos de registros mais informais ( como gírias, jargões, frases feitas, ditados populares, termos regionais ).
Assim, esperava -se que o participante : respeitasse a adequação vocabular, não empregando gírias ou expressões coloquiais ; apresentasse repertório vocabular diversificado e preciso ; evitasse repetição desnecessária de palavras ; utilizasse um vocabulário mais formal, como solicitado por um texto dissertativo.
Essas quatro competências foram reunidas em três itens de avaliação, atendendo aos aspectos explicitados anteriormente : I - Aspectos Ortográficos ( 20 \% ) : domínio das convenções ortográficas : grafia de vogais e consoantes, uso de maiúsculas e minúsculas, emprego do hífen e da acentuação gráfica, segmentação de palavras ; II - Aspectos Textuais ( 40 \% ) : domínio dos procedimentos de estruturação textual : organização interna dos períodos, emprego de conectores para a articulação lógica entre os períodos e entre os parágrafos, emprego de marcas de referenciação lexical e pronominal ; utilização dos sinais de pontuação que contribuem para a organização lógica da frase ; III - Aspectos Morfossintáticos e Vocabulares ( 40 \% ) : domínio das regras de caráter morfossintático estabelecidas como modelares do ponto de vista da modalidade escrita formal da norma-padrão da Língua Portuguesa : concordância nominal e verbal, regência nominal e verbal, colocação pronominal, flexão nominal e verbal, correlação entre tempos e modos verbais, ausência de marcas de oralidade.
Outra competência combinada a esta na análise das respostas é a seleção vocabular adequada à modalidade escrita formal da Língua Portuguesa, tendo em vista a intersecção entre as duas do ponto de vista das exigências do registro formal da modalidade escrita da norma-padrão.
Análise das respostas dadas pelos estudantes A Questão 1 propiciou a oportunidade de maior desenvolvimento do tema solicitado, gerando consequências na estruturação textual das respostas, já que os textos foram mais opinativos, mais elaborados e fluentes.
Já a Questão 2 teve um comportamento distinto, devido ao seu caráter mais expositivo, os textos foram mais curtos e o vocabulário mais limitado.
Observaram -se, nas duas questões, alguns casos de comprometimento de autoria, com textos muito próximos a paráfrases dos textos motivadores e sem marcas de repertório sociocultural produtivo.
Apresentam -se, a seguir, algumas \textbf{observações} sobre o desempenho dos participantes : I - Aspectos Ortográficos : o desempenho dos participantes revelou uma diferença muito grande nos dois aspectos analisados, como já foi observado nos relatórios dos anos anteriores : baixo índice de desvios da grafia padrão e grande índice de desvios de acentuação.
Em vários casos, ocorre ausência quase completa de acentuação gráfica.
Vale ressaltar que muitos participantes apresentaram uma caligrafia quase ininteligível.
Ainda nesse aspecto, ocorreu a omissão inadequada de cedilha, de til e do traço do T em algumas respostas.
Houve, inclusive, respostas com ausência total de cedilha ou til.
Os casos mais sistemáticos de eliminação do acento indicador da sílaba tônica foram : palavras proparoxítonas ( “ artistica ”, “ ludico ”, “ publico ”,"musica") ; palavras paroxítonas terminadas em ditongo crescente ( “ necessario ”,"historia","propria", “ vivencias ” ) ; palavras oxítonas ( “ tambem ”,"populacao").
Por outro lado, destaca -se o uso indevido do acento gráfico em palavras como “ propíciar ” “ criânça ”, por exemplo.
Quanto aos desvios de grafia, observam -se os seguintes exemplos : emprego inadequado de maiúscula : “ a Liberdade artística ” ; ou de minúscula :"constituição"; confusão entre"ss","sc","c"e"ç": “ condenassão ”, “ nessessario", “ repreiensão ”, “ conciencia ” ; eliminação de grafemas : “ essescia","questõe"; - desvios de segmentação : “ ser vidor", “ igual dade ”, “ possi bilidades"; troca de grafema : “ convença ” ( por “ conversa ” ), “ comun ”.
II - Aspectos Textuais : a estruturação textual foi a competência que apresentou o maior número de desvios, sendo responsável, em alguns textos, pela falta de clareza e pela desorganização.
No conjunto dos textos analisados, observaram -se diferentes níveis de domínio dos recursos de estruturação e organização textual.
No entanto, os participantes podem ser agrupados em três faixas descritas a seguir : - os que possuem grande domínio dos recursos de estruturação textual, com produção de textos que atendem às exigências de clareza e coerência, com progressão temática bem direcionada e utilização produtiva de operadores argumentativos ; - os que apresentam um domínio regular dos recursos de estruturação textual, com desvios na organização do texto e ausência de paragrafação, além de uso insuficiente de operadores argumentativos e processos de encaixamento sintático ; - os que apresentam baixo domínio dos recursos de estruturação textual, com muitos desvios que retomam hábitos da oralidade e que comprometem a coerência do texto : ocorrência de sequência justaposta de ideias sem encaixamentos sintáticos ; redução de estruturas subordinadas, ao lado do aumento na frequência de estruturas coordenadas e absolutas ; redução no uso de conectores para expressar relações lógicas essenciais à construção do texto ; emprego equivocado de operadores que não estabelecem relações lógicas coerentes entre ideias do texto ; fragmentação frasal que compromete a estrutura lógico-gramatical.
Uma estrutura que vem se ampliando na escrita dos participantes, refletindo hábitos da oralidade, é a duplicação de um dos elementos da frase, ou seja, a retomada anafórica de um elemento topicalizado, como por exemplo :"Mesmo com a época da censura, o ser humano ELE pode se expressar de várias formas..."; “ A cultura ELA vem trazida de várias épocas... ” ; “ Arte e cultura, ELAS estão relacionadas em uma só ” ; “ A cultura ELA vem enraizada a pessoa... ” ; “ As ações educativas ELAS podem ser consideradas boas ou ruins. ” Também em relação ao emprego dos sinais de pontuação, observaram -se distintos tipos de desempenho : participantes que dominam a funcionalidade da pontuação para a construção de textos fluentes e coesos, por um lado, e um grande número de participantes que apresentam precariedade em sua utilização, chegando ao ponto de elaborar parágrafos sem qualquer marca interna de pontuação para separar os períodos.
São os seguintes os tipos de desvios encontrados : - emprego de vírgula para separar o sujeito e o predicado : “ Os teatros, festas de rua, carnaval, elevam nosso conhecimento ” ( no lugar de “ Os teatros, festas de rua, carnaval elevam nosso conhecimento ” ) ;"que governos e municípios, façam a inclusão ” ( no lugar de “ que governos e municipios façam a inclusão ” ) ; - ocorrência de apenas uma das vírgulas para separar uma palavra, uma expressão ou uma oração encaixada : “ A liberdade artística que está presente em nossa Constituição 88, nos calça ” ( no lugar de “ A liberdade artística, que está presente em nossa Constituição 88, nos \textbf{calça} ) ” ; - emprego de ponto no lugar de vírgula para separar orações de um mesmo período : “ criando debate em espaços públicos.
Além de convidar os críticos a mostrarem o que é arte! ” ( no lugar de “ criando debate em espaços públicos, além de convidar os críticos a mostrarem o que é arte! ” ) ; - ausência de vírgula ou conjunção para separar elementos de uma enumeração :"As cidades inteligentes são caracterizadas por inteligência artificial, desenvolvimento sustentável conservação dos espaços verdes"( no lugar de"As cidades inteligentes são caracterizadas por inteligência artificial, desenvolvimento sustentável e conservação dos espaços verdes") ; - ausência de vírgula para separar oração adjetiva explicativa ou utilização inadequada para separar oração adjetiva restritiva :"As cidades inteligentes que são lugares em que a tecnologia avançada é utilizada..."( no lugar de"As cidades inteligentes, que são lugares em que a tecnologia avançada é utilizada,...") ; - emprego indevido do sinal de dois pontos : “ percebe -se um círculo onde estão infinitamente ligadas : a cultura e a arte ” ( no lugar de “ percebe -se um círculo onde estão infinitamente ligadas a cultura e a arte ” ) ; - ausência de pronome no processo de referenciação : “ censura a arte ou até mesmo impulsiona ” ( no lugar de “ censura a arte ou até mesmo a impulsiona ” ) ; - emprego de vírgula em lugar do ponto-final : “ são muitos os agravamentos em relação à tal situação, ” III - Aspectos Morfossintáticos e Vocabulares : o domínio das convenções morfossintáticas apresentou resultados que evidenciam, claramente, alguns aspectos problemáticos no desempenho dos participantes.
Assim como na análise dos desvios textuais, pode -se observar uma distinção entre três grupos de participantes : o grupo que revela domínio das convenções morfossintáticas ; o grupo que apresenta um domínio regular, com desvios pontuais ; e o grupo que revela domínio precário dessas convenções.
Pode -se destacar a regência, tanto verbal quanto nominal, como o aspecto com maior número de desvios : a ) Regência nominal e verbal : Na maioria dos textos, destacam -se os desvios relacionados à noção de"regência", ou seja, à relação de subordinação que se estabelece entre as palavras por meio, ou não, de preposições, refletindo uma tendência na modalidade escrita da língua.
O desvio mais frequente foi a falta do sinal indicativo da crase, em contextos em que seria obrigatória, ou, ao contrário, um uso indevido em contextos em que o princípio não se aplica.
Os dois tipos de desvios revelam que o usuário não tem consciência de que, sob a forma do termo “ a ”, se pressupõe a presença de uma preposição “ a ”, exigida pela regência do termo anterior.
Exemplos : “ que se contrapõem a essência das artes ” ( no lugar de “ que se contrapõem à essência das artes. ” ) ; “ Esse ato pode levar a naturalização"( no lugar de “ Esse ato pode levar à naturalização") ; “ Além de convidar os críticos a mostrarem o que é arte à fazendo! ” ( no lugar de “ Além de convidar os críticos a mostrarem o que é arte a fazendo! ” ).
Outro desvio relacionado à regência verbal e à nominal, encontrado frequentemente nas questões, é a ausência de preposição antes de pronome relativo, refletindo um processo generalizado na modalidade oral da língua, em situações de registro informal.
Exemplo :"a cidade que todos querem morar é..."( no lugar de"a cidade em que todos querem morar é...") b ) Concordância nominal e verbal : Os desvios relativos à aplicação das regras de concordância verbal e nominal correspondem a fatos que têm sido amplamente estudados na modalidade oral da Língua Portuguesa.
Quanto à concordância de número, observou -se ausência de marca ( com sujeito anteposto ou posposto ), como por exemplo a ausência de acento circunflexo na forma plural do presente do indicativo do verbo “ ter ”, considerada como um desvio na concordância verbal e não na acentuação gráfica.
Exemplo :"os governantes tem procurado planejar"( no lugar de"os governantes têm procurado planejar") ; “ onde é mostrado os acontecimentos da atualidade ” ( no lugar de “ onde são mostrados os acontecimentos da atualidade ” ).
Outro desvio na aplicação da regra da concordância de número foi o aparecimento da marca de plural em verbos ou adjetivos relacionados a núcleos substantivos no singular, devido à presença de modificadores no plural, evidenciando um processo de hipercorreção.
Esse processo tem se ampliado nas situações de oralidade, caracterizando uma tendência de variação da Língua Portuguesa.
Exemplos :"a construção de teatros e cinemas revelam uma preocupação das cidades inteligentes"( no lugar de"a construção de teatros e cinemas revela uma preocupação das cidades inteligentes") ;"o aumento dos acidentes ambientais devem ser evitados ” ( no lugar de"o aumento dos acidentes ambientais deve ser evitado ” ).
Quanto à concordância de gênero, vários casos foram observados, normalmente no âmbito de sintagmas nominais longos, em que o adjetivo está afastado do substantivo, conforme apontam as pesquisas sobre a modalidade oral.
Exemplos : “ sendo valorizado uma normatização de discurso que se contrapõe ” ( no lugar de “ sendo valorizada uma normatização de discurso que se contrapõe ” ) ; “ podendo essa agradar um público e desagradar outra ” ( no lugar de “ podendo essa agradar um público e desagradar outro ” ). c ) Registro da forma verbal “ é ” sem acento, provocando ambiguidade na leitura pela confluência com a conjunção aditiva"e".
Exemplo : “ A principal forma de realizar essa disseminação e criar nas escolas... ” ( no lugar de “ A principal forma de realizar essa disseminação é criar nas escolas... ” ). d ) Emprego do pronome relativo “ onde ” como relativo universal, típica marca da oralidade, em que esse pronome não é mais utilizado apenas em sua função locativa"lugar onde se realiza a ação".
Exemplo :"a aceitação não é mais a mesma, onde muitas vezes é a cultura que censura a arte"; “ situações onde ” ; “ cultura mais cotidiana, onde ”. e ) Falta de artigo definido antes de substantivo, comprometendo o paralelismo em enumerações : Exemplo :"censurar a arte e cultura é prejudicial"( no lugar de"censurar a arte e a cultura é prejudicial"). f ) Desvios relacionados à flexão verbal : Exemplos : emprego indevido do tempo ou do modo verbal : “ adaptado ” por “ adaptando ” ; eliminação de marca de infinitivo : “ ta ” por “ estar ”, “ está ” por “ estar ”, “ esta ” por “ estar ” ; “ estás ” por “ estas ” ; utilização de construção de “ estar ” + gerúndio como marca de oralidade : “ devemos cada vez mais ta melhorando ” ; segmentação indevida da desinência de número e pessoa de verbos : “ entender -mos ”, “ aprender -mos ”.
Assim como nas demais competências avaliadas, há diferenças entre os participantes em relação à seleção e ao emprego do vocabulário.
Foram encontrados textos com alto padrão de domínio vocabular, chegando até ao emprego de vocabulário técnico das áreas do saber abordadas nas questões : arte, cultura, meio ambiente e tecnologia.
Por outro lado, na maioria dos textos foram observados vários tipos de inadequação na seleção e na utilização do vocabulário, sendo que o principal deles foi a excessiva repetição de certas palavras, revelando limitação de repertório vocabular.
Além disso, observou -se a ocorrência de expressões da oralidade, revelando falta de domínio de vocabulário mais abstrato e de maior complexidade, essencial ao desenvolvimento do texto de base dissertativa.
Exemplos : seleção vocabular indevida : “ crescemos tão intrínsecos a esses usos... ” ; emprego inadequado do verbo “ ter ” no lugar de “ haver ” : “ Durante a ditadura teve muita repressão ” ; imprecisão vocabular : “ ela anda lado a lado com a existência do que conhecemos de cultura ”, “ A liberdade artística que está presente em nossa Constituição 88, nos \textbf{calça} e... ”.
A imprecisão vocabular, como o uso da palavra “ coisa ” em vários textos, demonstra o domínio precário da norma-padrão da Língua Portuguesa, com presença constante de traços de oralidade. 6.3.1.8 CONSIDERAÇÕES FINAIS SOBRE A CORREÇÃO DE LÍNGUA PORTUGUESA DAS RESPOSTAS ÀS QUESTÕES DISCURSIVAS DO COMPONENTE DE FORMAÇÃO GERAL O conjunto de textos analisados no Enade 2021 permite traçar um perfil bem detalhado do desempenho linguístico do estudante de graduação brasileiro.
Conforme ficou comprovado nos comentários dos exemplos de respostas apresentados neste relatório, uma parcela significativa dos participantes revela lacunas importantes em sua formação, principalmente no que diz respeito à estruturação e à organização de textos.
Como esta é uma competência essencial à produção escrita, que deve ser construída ao longo dos anos de escolaridade ( do ensino fundamental ao ensino superior ), sua recuperação é trabalhosa e depende de um esforço intenso de exposição a textos e de reflexão sobre os requisitos que garantem a clareza e o encadeamento das ideias. 6.3.2 COMPONENTE DE CONHECIMENTO ESPECÍFICO Na parte da prova relativa às Questões Discursivas do Componente de Conhecimento Específico ( Tabela 6.17a ), observa -se que a Média foi mais baixa do que para as Questões Discursivas do Componente de Formação Geral.
Enquanto, no Componente de Formação Geral, a Média para estudantes de Geografia ( Licenciatura ) de todo o Brasil foi 27,0, na parte de Conhecimento Específico, a Média foi 18,3.
A maior Média deste componente foi obtida pelos estudantes da região Sudeste ( 22,1 ), e a menor, pelos das regiões Norte e Sul ( 16,6 ).
As demais regiões obtiveram Média igual a 17,0 ( Nordeste ) e 17,1 ( Centro-Oeste ).
Quanto à variabilidade das notas, o Desvio padrão de todo o Brasil, entre os alunos de Licenciatura, foi A competência relacionada ao domínio do vocabulário, que deve ser diversificado e apropriado ao padrão formal da língua, também apresentou um desempenho insatisfatório pela maioria dos participantes, que se limitam ao senso comum, com vocabulário genérico e vago, o que é inadequado para atender ao perfil de um graduando.
As outras duas competências avaliadas ( convenções ortográficas, convenções morfossintáticas ) apresentam desvios pontuais e são de mais fácil recuperação, com exceção dos desvios morfossintáticos que já estão incorporados na modalidade oral informal da Língua Portuguesa, configurando processos de variação linguística e, talvez, antecipando futuras mudanças no padrão da modalidade escrita.
Com base na avaliação do desempenho linguístico aqui relatada, conclui -se que é preciso aperfeiçoar a formação superior, dando aos estudantes mais acesso às ferramentas necessárias para que sejam bem-sucedidos nas diferentes situações de sua vida cidadã – profissional, cultural e política.
Essa é uma contribuição efetiva que os resultados do Enade podem oferecer às instituições de ensino superior do país.
Para finalizar, vale observar que as duas questões discursivas atenderam às exigências de avaliação, tanto do domínio discursivo dos estudantes, quanto do seu conhecimento sobre temas relacionados ao conhecimento acadêmico, o que pode ser extremamente útil para que os cursos de ensino superior possam contribuir na formação de seus graduandos.
Se o objetivo é avaliar a competência discursiva dos formandos brasileiros, a elaboração de uma resposta de base dissertativo-argumentativa ( Questão 1 ) e de uma resposta de base expositiva ( Questão 2 ) certamente ofereceu material relevante para essa análise, por incentivar a reflexão e a organização do pensamento por meio de uma sequência textual coerente e complexa. 6.3.2 COMPONENTE DE CONHECIMENTO ESPECÍFICO Na parte da prova relativa às Questões Discursivas do Componente de Conhecimento Específico ( Tabela 6.17a ), observa -se que a Média foi mais baixa do que para as Questões Discursivas do Componente de Formação Geral.
Enquanto, no Componente de Formação Geral, a Média para estudantes de Geografia ( Licenciatura ) de todo o Brasil foi 27,0, na parte de Conhecimento Específico, a Média foi 18,3.
A maior Média deste componente foi obtida pelos estudantes da região Sudeste ( 22,1 ), e a menor, pelos das regiões Norte e Sul ( 16,6 ).
As demais regiões obtiveram Média igual a 17,0 ( Nordeste ) e 17,1 ( Centro-Oeste ).
Quanto à variabilidade das notas, o Desvio padrão de todo o Brasil, entre os alunos de Licenciatura, foi 16,2.
O maior Desvio padrão foi encontrado na região Sudeste ( 17,9 ), e o menor, na região Nordeste ( 14,6 ).
A maior nota Máxima, 93,3, foi obtida na região Sudeste.
Nas demais regiões, as notas Máximas foram : 71,7 ( Norte ), 76,7 ( Nordeste e Centro-Oeste ) e 83,3 ( Sul ).
A nota Mínima zero foi obtida por, pelo menos, um estudante de Licenciatura em todas as regiões do Brasil.
A Mediana do Brasil foi igual a 15,0, enquanto na região Sudeste foi 20,0 e nas demais regiões foi 13,3.
Na parte da prova relativa às Questões Discursivas do Componente de Conhecimento Específico ( Tabela 6.17b ), observa -se que a Média foi mais baixa do que para as Questões Discursivas do Componente de Formação Geral.
Enquanto, no Componente de Formação Geral, a Média para estudantes de Geografia ( Bacharelado ) de todo o Brasil foi 32,9, na parte de Conhecimento Específico, a Média foi 19,8.
A maior Média deste componente foi obtida pelos estudantes da região Sudeste ( 20,7 ), e a menor, pelos da região Centro-Oeste ( 16,0 ).
As demais regiões obtiveram Média igual a 17,6 ( Norte ), 19,6 ( Nordeste ) e 20,4 ( Sul ).
Quanto à variabilidade das notas, o Desvio padrão de todo o Brasil, entre os alunos de Bacharelado foi 15,3.
O maior Desvio padrão foi encontrado na região Nordeste ( 16,0 ), e o menor, na região Centro-Oeste ( 13,7 ).
A maior nota Máxima, 81,7, foi obtida na região Sudeste.
Nas demais regiões, as notas Máximas foram : 61,7 ( Norte ), 71,7 ( Nordeste ), 63,3 ( Sul ) e 60,0 ( Centro-Oeste ).
A nota Mínima zero foi obtida por, pelo menos, um estudante de Bacharelado em todas as regiões do Brasil.
A Mediana do Brasil e da região Sul foi igual a 18,3 ; na região Norte foi 13,3 ; na região Nordeste foi 16,7 ; na região Sudeste, foi 20,0 – a maior ; e na região Centro-Oeste, foi 15,0.
A maior nota Máxima, 81,7, foi obtida na região Sudeste.
Nas demais regiões, as notas Máximas foram : 61,7 ( Norte ), 71,7 ( Nordeste ), 63,3 ( Sul ) e 60,0 ( Centro-Oeste ).
A nota Mínima zero foi obtida por, pelo menos, um estudante de Bacharelado em todas as regiões do Brasil.
A Mediana do Brasil e da região Sul foi igual a 18,3 ; na região Norte foi 13,3 ; na região Nordeste foi 16,7 ; na região Sudeste, foi 20,0 – a maior ; e na região Centro-Oeste, foi 15,0.
Na região região Nordeste foi 16,7 ; na região Sudeste, foi 20,0 – a maior ; e na região Centro-Oeste, foi 15,0.
No Gráfico 6.10, representa -se a distribuição das notas nas Questões Discursivas no Componente de Conhecimento Específico.
A moda dessa distribuição ocorre no intervalo [ 0 ; 10 ], tanto para os alunos de Licenciatura quanto para os de Bacharelado, com 39,8 \% e 34,0 \% do total de participantes, respectivamente.
Observa -se que nesse intervalo estão computadas as ocorrências de respostas em branco para as três Questões Discursivas de Conhecimentos Específicos.
A análise de cada uma destas questões será feita a seguir. 6.3.2.1 ANÁLISE DA QUESTÃO DISCURSIVA 3 DO COMPONENTE DE CONHECIMENTO ESPECÍFICO A Questão Discursiva 3 de Conhecimentos Específicos foi a mesma para as duas habilitações.
Os resultados para a Licenciatura constam da Tabela 6.18a e para o Bacharelado estão apresentados na Tabela 6.18b.
Cabe esclarecer que por terem respondido à mesma questão os resultados são comparáveis.
Na Questão 3, cujos resultados aferidos se encontram descritos na Tabela 6.18a, a Média dos estudantes de Licenciatura de todo o Brasil foi 16,8.
O desempenho dos estudantes de Licenciatura nessa questão foi superior ao da Questão 4 e inferior ao da Questão 5.
A menor Média nessa questão foi obtida pelos estudantes de Licenciatura da região Norte ( 14,7 ), enquanto a maior Média foi obtida na região Sudeste ( 19,7 ).
Quanto à variabilidade das notas, o Desvio padrão dos estudantes de Licenciatura de todo o Brasil foi 17,8.
O maior Desvio padrão foi obtido na região Sudeste ( 19,2 ), enquanto o menor foi obtido na região Sul ( 16,7 ).
A Mediana ( 10,0 ), a nota Máxima ( 100,0 ) e a nota Mínima ( 0,0 ) foram iguais para todas as regiões, sendo a Mediana da região Sudeste ( 20,0 ) a única exceção.
As Médias para os alunos de Licenciatura ( ver Tabela 6.18a ) foram, respectivamente para os de cursos Presenciais e a Distância, 18,6 e 14,6 ( 16,8 para o total ).
A nota Máxima ( 100,0 ) foi atingida por, pelo menos, um concluinte em quase todas as regiões, com exceção da região Sul ( 90,0 ), para a Modalidade Presencial, e das regiões Nordeste ( 60,0 ) e Centro-Oeste ( 75,0 ) para a Modalidade a Distância.
A Mediana foi 20,0 para os cursos Presenciais e 10,0 para os a Distância.
Cumpre notar que somente na região Norte não houve participação de estudantes de cursos a Distância de Licenciatura.
A nota Mínima foi zero para todas as combinações de Modalidade de Ensino e região.
Na Questão 3, cujos resultados aferidos se encontram descritos na Tabela 6.18b, a Média dos estudantes de Bacharelado de todo o Brasil foi 15,1.
O desempenho dos estudantes de Bacharelado nessa questão foi inferior ao verificado nas Questões 4 e 5.
A menor Média nessa questão foi obtida pelos estudantes das regiões Norte e Centro-Oeste ( 10,9 ), enquanto a maior Média foi obtida na região Sul ( 16,4 ).
Quanto à variabilidade das notas, o Desvio padrão dos estudantes de Bacharelado de todo o Brasil foi 16,0.
O maior Desvio padrão foi obtido na região Nordeste ( 16,6 ), enquanto o menor foi obtido na região Centro-Oeste ( 13,9 ).
A nota Máxima, 50,0 pontos, foi alcançada por, pelo menos, um estudante de Bacharelado em todas as regiões.
A Mediana dos estudantes de Bacharelado de todo o Brasil, foi 10,0, a mesma para quase todas as regiões, com exceção das regiões Norte e Centro-Oeste, que obtiveram Mediana igual a zero, indicando que, nessas regiões, pelo menos 50 \% dos estudantes tiraram nota zero.
A nota Mínima ( 0,0 ) foi a mesma em todas as regiões do Brasil.
As Médias para os alunos de Bacharelado ( ver Tabela 6.18b ) foram, respectivamente para os de cursos Presenciais e a Distância, 15,0 e 16,0 ( 15,1 para o total ).
A nota Máxima foi 50,0 para todas as combinações de Modalidade de Ensino e região.
A Mediana foi 10,0 para os cursos Presenciais e 17,5 para os a Distância.
Cumpre notar que somente nas regiões Sudeste e Sul houve participação de estudantes de cursos a Distância de Bacharelado.
A nota Mínima foi zero para todas as combinações de Modalidade de Ensino e região.
No Gráfico 6.11, mostra -se a distribuição das notas na Questão Discursiva 3, do Componente de Conhecimento Específico da Área de Geografia ( Licenciatura/Bacharelado ).
Sem levar em consideração que 26,2 \% dos estudantes de Licenciatura tiraram nota zero e 11,5 \% deixaram a resposta à questão em branco perfazendo um total de 37,7 \% o contingente daqueles que zeraram a questão, a moda dessa distribuição seria o intervalo ( 20 ; 30 ], com 20,1 \% do total de participantes de Licenciatura.
Dentre os alunos de Bacharelado, sem levar em consideração que 33,4 \% tiraram nota zero e 10,1 \% deixaram a resposta à questão em branco perfazendo um total de 43,5 \% o contingente que zerou a questão, a moda dessa distribuição seria o intervalo ( 20 ; 30 ], com 31,5 \% do total de participantes. 6.3.2.2 COMENTÁRIOS SOBRE AS RESPOSTAS À QUESTÃO DISCURSIVA 3 O enunciado da Questão 3, que incluía uma coleção de mapas, era claro e abordava conteúdo adequado para avaliar estudantes concluintes do curso de Geografia do Bacharelado e da Licenciatura.
A questão pode ser considerada com nível de profundidade mediano.
O item ‘ a ’ solicitava que o respondente analisasse"evolução"do fluxo de passageiros e sua"distribuição".
Assim, deveria verificar não somente as distintas temporalidades da evolução do fluxo de passageiros observando os dados dos mapas, mas também, deveria apontar os fatores que levaram a essa evolução e a distribuição pelas diferentes regiões brasileiras.
Nas respostas a este item, parece ter havido alguma confusão de interpretação por parte dos respondentes já que"evolução"do fluxo de passageiros e sua"distribuição"são aspectos inter-relacionáveis, levando alguns estudantes a tratarem estas duas perspectivas de análise de forma trocada.
No item ‘ b ’, o respondente deveria propor uma atividade didática correlacionando"transporte aéreo"e"desigualdade social no Brasil", o que, em função do caráter abrangente do comando, permitia diversas possibilidades de respostas, além das duas exemplificadas pelo padrão de resposta.
No que diz respeito às competências e habilidades presentes nas Diretrizes Curriculares da Área, a questão ( itens ‘ a ’ e ‘ b ’ ) avaliava tópicos presentes no Art. 7º da Portaria nº 378 ( Diretrizes da Área de Geografia/Licenciatura ).
Entretanto, em relação ao Bacharelado, o item ‘ b ’, que solicitava a apresentação de uma atividade didática para estudantes do sétimo ano do ensino fundamental, não era adequado, e os conteúdos necessários não são contemplados no artigo 6º da Portaria nº 379, que trata das competências dos formandos em Bacharelado.
Por isso, este item, para os estudantes de Bacharelado restou anulado pela Comissão Assessora da Área.
Uma parte das respostas dadas pelos estudantes, particularmente para o item ‘ a ’, foi contemplada pelo padrão de respostas.
Entretanto, muitos não separaram, no texto redigido, a evolução do fluxo de passageiros, de sua distribuição pelas áreas geográficas.
O foco da maioria das respostas ficou limitado a comentar as facilidades oferecidas na aquisição de passagens, seja pela melhora do poder aquisitivo da população, seja pela possibilidade de compra de passagens aéreas usando parcelamentos e/ou promoções.
A melhoria na infraestrutura de aeroportos, em cidades de pequeno e médio portes, também foi bastante citada.
O item ‘ b ’ foi equivocadamente respondido por muitos, isso porque exigia conhecimentos de práticas pedagógicas em Geografia que a grande maioria dos licenciandos demonstrou desconhecer.
Neste caso, poucos se aproximaram do padrão de resposta.
Foram dadas muitas respostas com sugestões de levantamento de dados pelos alunos, como por exemplo, estudantes do 7º ano do ensino fundamental deveriam entrevistar familiares e amigos, perguntando se eles já tinham realizado viagens aéreas e o poder aquisitivo de cada entrevistado.
No geral, muitos concluintes responderam apenas ao item ‘ a ’, em linguagem compatível com os termos usualmente aplicados em Geografia.
Apesar disso, muitas \textbf{repostas} se limitavam à interpretação dos mapas apresentados no enunciado da questão.
Nessas respostas, alguns conseguiram discorrer sobre a evolução e o fator de distribuição de passageiros pelas regiões brasileiras, mas sem distinguir um processo do outro, como já comentado.
Este foi o principal erro cometido.
Com relação ao item ‘ b ’, não houve diversidade significativa de respostas, já que uma boa parcela daqueles que tentaram responder a esse item, como já comentado, sugeria atividades nas quais os estudantes deveriam realizar pesquisas de vários tipos ( de preços de passagens aéreas e comparação com o salário mínimo ; aplicar questionário em casa para saber com que frequência as pessoas utilizam transporte aéreo e avaliar com o poder aquisitivo ), sem propor ações pedagógicas, análise e interpretação de resultados ou discussão sobre desigualdades sociais, por exemplo.
O desempenho dos estudantes, nessa questão, foi ruim ( predomínio de notas muito fracas ).
Isso pode retratar a pouca capacidade de entendimento do que foi solicitado, em ambos os itens formulados, com o agravante de o item ‘ b ’ ter sido respondido por pouquíssimos estudantes da Licenciatura.
Não houve diferenciações significativas entre as respostas de estudantes do Bacharelado e da Licenciatura em relação ao item ‘ a ’, visto que o item ‘ b ’ foi anulado para o Bacharelado.
A maioria das respostas foi muito fraca, retratando uma grande deficiência no ordenamento das ideias nas redações, especialmente para o item ‘ b ’ da Licenciatura.
Apenas um estudante da Licenciatura obteve nota máxima.
A preocupação central identificada nas respostas redigidas pelos estudantes esteve voltada, principalmente, para mostrar o baixo poder aquisitivo das pessoas tornando difícil adquirir uma passagem aérea, em detrimento da análise do quadro evolutivo e da distribuição, pelo país, do modal aéreo ( item ‘ a ’ ).
Isso revelou graves deficiências de \textbf{aprendizado} sobre assuntos geográficos de cunho geral. 6.3.2.3 ANÁLISE DA QUESTÃO DISCURSIVA 4 DO COMPONENTE DE CONHECIMENTO ESPECÍFICO A Questão Discursiva 4 de Conhecimentos Específicos foi a mesma para as duas Habilitações.
Os resultados para a Licenciatura constam da Tabela 6.19a e para o Bacharelado estão apresentados na Tabela 6.19b.
Cabe esclarecer que por terem respondido à mesma questão os resultados são comparáveis.
Na Tabela 6.19a, constam as informações relativas à Questão 4 do conjunto de Questões do Componente de Conhecimento Específico.
O desempenho dos estudantes de Licenciatura nessa questão foi inferior ao das Questões 3 e 5.
A Média geral dos estudantes de Licenciatura de todo o Brasil foi 15,8, sendo a maior Média registrada na região Norte ( 19,2 ), e a menor, na região Centro-Oeste ( 12,5 ).
A nota Máxima ( 100,0 ) foi atingida por, pelo menos, um concluinte de Licenciatura em todas as regiões.
A Mediana dos estudantes de Licenciatura em todo o Brasil foi zero, o mesmo valor foi obtido nas regiões Nordeste, Sul e Centro-Oeste, indicando que, no Brasil como um todo e nessas regiões, pelo menos 50 \% dos estudantes tiraram nota zero.
As demais regiões obtiveram Mediana igual a 10,0.
Foi zero a nota Mínima de todas as regiões.
As Médias para os alunos de Licenciatura ( ver Tabela 6.19a ) foram, respectivamente para os de cursos Presenciais e a Distância, 17,9 e 13,0 ( 15,8 para o total ).
A nota Máxima ( 100,0 ) e a nota Mínima ( 0,0 ) foram as mesmas em todas as combinações de Modalidade de Ensino e região.
A Mediana foi zero para os cursos a Distância em todas as regiões e 10,0 para o total dos cursos Presenciais.
Cumpre notar que somente na região Norte não houve participação de estudantes de cursos a Distância de Licenciatura.
Na Tabela 6.19b, constam as informações relativas à Questão 4 do conjunto de Questões do Componente de Conhecimento Específico.
O desempenho dos estudantes de Bacharelado nessa questão foi superior ao verificado nas Questões 3 e 5.
A Média geral dos estudantes de Bacharelado do Brasil foi 28,1, sendo a maior Média registrada na região Sudeste ( 30,1 ), e a menor, nas regiões Norte e Centro-Oeste ( 24,6 ).
A nota Máxima ( 100,0 ) foi atingida por, pelo menos, um concluinte de Bacharelado de cada região.
A Mediana em todo o Brasil foi 20,0, mesmo valor foi obtido em quatro das cinco Grandes Regiões, com exceção da região Norte ( 10,0 ).
A nota Mínima de todas as regiões foi igual a zero.
As Médias para os alunos de Bacharelado ( ver Tabela 6.19b ) foram, respectivamente para os de cursos Presenciais e a Distância, 28,5 e 19,5 ( 28,1 para o total ).
A nota Máxima ( 100,0 ), a nota Mínima ( 0,0 ) e a Mediana ( 20,0 ) foram as mesmas quase em todas as combinações de Modalidade de Ensino e região, com exceção das regiões Sudeste e Sul houve participação de estudantes de cursos a Distância de Bacharelado.
No Gráfico 6.12, está representada a distribuição de notas da Questão Discursiva 4, do Componente de Conhecimento Específico.
Sem levar em consideração que 15,8 \% dos estudantes de Licenciatura deixaram a resposta à questão em branco, e que 37,1 \% tiraram nota zero, mesmo respondendo essa questão, a moda dessa distribuição seria o intervalo ( 10 ; 20 ], com 15,5 \% do total de participantes de Licenciatura.
Dentre os alunos de Bacharelado, sem levar em consideração que 15,1 \% deixaram a resposta à questão em branco, e 19,1 \% receberam nota 0,0, a moda dessa distribuição seria o intervalo ( 10 ; 20 ], com 16,7 \% do total de participantes. 6.3.2.4 COMENTÁRIOS SOBRE AS RESPOSTAS À QUESTÃO DISCURSIVA 4 O enunciado da Questão 4, aplicada tanto para concluintes do Bacharelado quanto da Licenciatura, continha um texto motivador claro e que contribuía para a construção das respostas.
Os comandos dos itens ‘ a ’ e ‘ b ’ eram diretos e de fácil compreensão.
O item ‘ a ’ exigia respostas objetivas, sem possibilidades de variações.
Já o item ‘ b ’, que solicitava a apresentação de três vantagens de conciliar a proteção da natureza e a associabilidade com povos tradicionais, possibilitava grande variedade de respostas.
Muitas das respostas possíveis foram listadas no padrão de respostas.
A questão avaliava conteúdos, competências e habilidades constantes das Diretrizes Curriculares, tanto do Bacharelado, quanto da Licenciatura, tratando de temáticas ambientais e socioculturais atuais, amplamente divulgadas nos diversos meios de comunicação.
Neste contexto, o nível de profundidade exigido foi de baixo a mediano, sendo o item ‘ b ’ o mais complexo, já que envolvia explicações sobre cada vantagem mencionada.
Dentre as três questões discursivas da Área, essa questão foi a mais fácil de ser respondida.
O padrão de resposta de ambos os itens atendia plenamente ao que se esperava encontrar como possíveis respostas.
Entretanto, muitos estudantes discorreram sobre a questão ambiental brasileira sem se reportar diretamente ao que foi solicitado.
Nas respostas ao item ‘ a ’, muitos estudantes contextualizaram a questão ambiental atual do Brasil, destacando a necessidade de proteção urgente do que ainda resta de ecossistemas, sem discriminar a diferença entre conservar e preservar.
Vários respondentes inverteram os dois conceitos ( de preservação e de conservação ), o que elevou significativamente o número de erros nesse item da questão.
Já no item ‘ b ’, frequentemente, as respostas destacavam a necessidade das populações ( tradicionais ou não ) protegerem o meio ambiente, fazendo o seu papel de exercício da cidadania e/ou auxiliando o poder público na proteção dos recursos naturais.
Poucos foram os que discriminaram, com clareza, as três vantagens solicitadas.
Dos que citaram pelo menos uma vantagem, muitos destacaram a importância das populações tradicionais na proteção ambiental, mas sem dizer de que maneira ocorre essa atuação/relevância.
Foi também significativo o número dos que discorreram sobre a importância de assegurar a conservação do meio ambiente através da inclusão social, sem, necessariamente, responder ao que foi solicitado.
Apesar destes problemas, o aproveitamento das respostas desse item foi mais significativo.
O desempenho dos estudantes, tanto do Bacharelado quanto da Licenciatura, foi fraco, apesar de o tema abordado pela questão ser de amplo conhecimento social e de relativa facilidade de compreensão.
Observou -se diferença no desempenho dos concluintes de Bacharelado em relação aos de Licenciatura.
Os bacharelandos responderam com mais propriedade aos dois itens, principalmente ao item ‘ b ’, apresentando respostas mais diretas e corretas, não parafraseando o enunciado, como ocorreu em muitas respostas da Licenciatura.
As respostas dadas pelos estudantes revelaram que, apesar da ampla divulgação de conhecimentos sobre as questões socioambientais brasileiras, quando se perguntam aspectos específicos, com uso de terminologias também mais específicas, como"conservação ” e “ preservação", muitos respondentes não conseguem especificar a diferença.
Demonstraram reconhecer a importância da manutenção dos valores e culturas das populações tradicionais, mas muitos não conseguiram apontar as vantagens de conciliá -los com a proteção da natureza.
Duas das questões objetivas poderiam contribuir para a elaboração de respostas para a Questão Discursiva 4.
A Questão Objetiva 13 tratava da proteção da identidade de povos tradicionais, ao abordar a ideia de"lugar"e “ comunidade", e a Questão Objetiva 25 abordava o tema da conservação ambiental proporcionada pelos Geoparques, principalmente do meio abiótico. 6.3.2.5 ANÁLISE DA QUESTÃO DISCURSIVA 5 DO COMPONENTE DE CONHECIMENTO ESPECÍFICO A Questão Discursiva 5 de Conhecimentos Específicos aplicada aos concluintes de Licenciatura foi diferente da Questão 5 respondida pelos concluintes do Bacharelado, portanto, por não terem respondido à mesma questão os resultados não são comparáveis.
As Tabelas 6.20a e 6.20b contém, respectivamente para Licenciatura e Bacharelado, as informações relativas à Questão 5 do Componente de Conhecimento Específico.
O desempenho dos estudantes de todo o Brasil nesta questão foi superior ao verificado nas Questões 3 e 4, desse componente.
A nota Média dos estudantes de todo o Brasil foi 22,4.
A maior Média foi registrada na região Sudeste ( 27,5 ), enquanto a menor Média foi registrada na região Norte ( 15,9 ).
Quanto à variabilidade das notas, o Desvio padrão dos estudantes do Brasil, como um todo, foi 25,7.
Enquanto o maior desvio foi encontrado na região Sudeste ( 28,4 ), o menor foi encontrado na região Norte ( 23,1 ).
A Mediana para o Brasil e em quase todas as regiões foi 15,0, com exceção da região Norte, onde a Mediana foi zero, indicando que, nessa região, pelo menos 50 \% dos estudantes tiraram nota zero.
Para o conjunto de estudantes de Geografia do Brasil, a nota Máxima foi 100,0, obtida por, pelo menos, um estudante das regiões Nordeste, Sudeste e Centro-Oeste, enquanto nas demais regiões tais valores foram 90,0.
A nota Mínima do Brasil foi zero, o mesmo valor encontrado em todas as regiões.
As Médias para os alunos de Licenciatura ( ver Tabela 6.20a ) foram, respectivamente para os de cursos Presenciais e a Distância, 23,7 e 20,7 ( 22,4 para o total ).
A nota Máxima ( 100,0 ) foi atingida por, pelo menos, um concluinte das regiões Nordeste, Sudeste e Centro-Oeste para a Modalidade Presencial, e da região Sudeste para a Modalidade a Distância.
A Mediana foi 15,0 para os cursos de ambas as Modalidades em todas as regiões, com exceção dos cursos Presenciais da região Norte e dos cursos a Distância das regiões Nordeste e Centro-Oeste, cujas Medianas foram zero, indicando que, nessas respectivas regiões e Modalidades, pelo menos 50 \% dos estudantes tiraram nota zero.
Cumpre notar que somente na região Norte não houve participação de estudantes de cursos a Distância de Licenciatura.
A nota Mínima foi zero para todas as combinações de Modalidade de Ensino e região.
Na Tabela 6.20b, constam as informações relativas à Questão 5 do conjunto do Componente de Conhecimento Específico.
O desempenho dos estudantes de todo o Brasil nesta questão foi superior ao verificado na Questão 3 e inferior ao verificado na Questão 5.
A nota Média dos estudantes de todo o Brasil foi 16,2.
A maior Média foi registrada na região Sul ( 17,7 ), enquanto a menor Média foi registrada na região Centro-Oeste ( 12,4 ).
Quanto à variabilidade das notas, o Desvio padrão dos estudantes do Brasil, como um todo, foi 18,7.
Enquanto o maior desvio foi encontrado na região Nordeste ( 20,0 ), o menor foi encontrado na região Centro-Oeste ( 15,4 ).
A Mediana para o Brasil foi 20,0, assim como as Medianas das regiões Norte, Sudeste e Sul.
Na região Nordeste a Mediana foi 10,0 e na região Centro-Oeste foi zero, indicando que, nessa região, pelo menos 50 \% dos estudantes tiraram nota zero.
Para o conjunto de estudantes de Geografia do Brasil, a nota Máxima foi 100,0, obtida por, pelo menos, um estudante das regiões Nordeste e Sudeste.
A nota Mínima do Brasil foi zero, o mesmo valor encontrado em todas as regiões.
As Médias para os alunos de Bacharelado ( ver Tabela 6.20b ) foram, respectivamente para os de cursos Presenciais e a Distância, 16,5 e 11,7 ( 16,2 para o total ).
A nota Máxima ( 100,0 ) foi atingida por, pelo menos, um concluinte das regiões Nordeste e Sudeste para a Modalidade Presencial, enquanto a nota Máxima dos cursos a Distância foi obtida na região Sudeste ( 60,0 ).
Nos cursos Presenciais a Mediana foi 20,0 nas regiões Norte, Sudeste e Sul, enquanto foi 10,0 para a região Nordeste e zero para a região Centro-Oeste, indicando que, para nos cursos Presenciais dessa região, pelo menos 50 \% dos estudantes tiraram nota zero para zero.
Para os cursos a Distância a Mediana foi zero para todas as regiões.
Cumpre notar que somente nas regiões Sudeste e Sul houve participação de estudantes de cursos a Distância de Bacharelado.
A nota Mínima foi zero para todas as combinações de Modalidade de Ensino e região.
No Gráfico 6.13, é representada a distribuição das notas da Questão Discursiva 5 do Componente de Conhecimento Específico.
Sem levar em consideração que 18,8 \% dos estudantes de Licenciatura deixaram a resposta à questão em branco, e 22,3 \% receberam nota 0,0, a moda dessa distribuição seria o intervalo ( 10 ; 20 ], com 20,7 \% do total de participantes de Licenciatura.
Dentre os alunos de Bacharelado, sem levar em consideração que 18,3 \% deixaram a resposta à questão em branco, e 27,5 \% receberam nota 0,0, a moda dessa distribuição seria o intervalo ( 10 ; 20 ], com 26,2 \% do total de participantes. 6.3.2.6 COMENTÁRIOS SOBRE AS RESPOSTAS À QUESTÃO DISCURSIVA 5 Questão 5 – Licenciatura A Questão 5, aplicada apenas a concluintes da Licenciatura, trazia um texto motivador que abordava o Dia Mundial do meio Ambiente em época de pandemia causada pelo Coronavírus.
O texto fornecia informações atuais, inclusive por levantar discussões sobre associação entre o meio ambiente e a saúde do planeta, a fome, a sustentabilidade, o que permitia ao estudante refletir sobre o que seria solicitado, a seguir.
A questão exigia do respondente conhecimentos sólidos sobre as mudanças climáticas para embasar boas respostas acerca da problemática socioambiental, o que inclui os efeitos \textbf{nocivos} das mudanças climáticas e o desafio da produção de alimentos.
No item ’ a ’, o respondente deveria apresentar e descrever duas ações sustentáveis voltadas à produção de alimentos para um maior número de pessoas.
Além disso, as ações apresentadas deveriam ser capazes de diminuir os efeitos das mudanças climáticas.
Como resposta ao item ‘ b ’, também eram solicitadas duas ações, mas, neste caso, o comando era “ apresente ” e não “ descreva ” como no item ‘ a ’.
Neste item, o grau de subjetividade era mais alto, já que não foi definido um objetivo claro, como no item ‘ a ’, para o qual as ações deveriam ser propostas.
Foram solicitadas ações que pudessem ser desenvolvidas “ pelos diversos segmentos sociais ” com vistas a uma sociedade sustentável, o que comporta respostas bastante variadas. subjetividade era mais alto, já que não foi definido um objetivo claro, como no item ‘ a ’, para o qual as ações deveriam ser propostas.
Foram solicitadas ações que pudessem ser desenvolvidas “ pelos diversos segmentos sociais ” com vistas a uma sociedade sustentável, o que comporta respostas bastante variadas.
O padrão de resposta, para os dois itens, listou algumas das principais linhas de ações que poderiam ser apresentadas, sem esgotar, obviamente, todas as respostas possíveis.
Vale destacar que o padrão de resposta proposto conseguiu prever uma boa variedade das respostas satisfatórias apresentadas pelos estudantes, principalmente para o item ‘ b ’.
Considerando -se que os estudantes deveriam apresentar quatro ações, sendo necessário descrever com mais detalhes as relativas ao item ‘ a ’, o espaço para resposta ( 15 linhas ) deveria ser maior.
Dentre os estudantes que apresentaram boas respostas, observou -se uma grande necessidade de síntese e o uso total do espaço disponível.
Uma parcela significativa dos estudantes não separou as quatro respostas em dois itens.
Foi comum apresentarem um texto único, às vezes sem quebra de parágrafos, que, de forma mesclada, buscava dar conta do que foi solicitado em cada item.
Foi muito frequente discorrerem sobre a substituição das monoculturas voltadas à exportação por sistemas agroflorestais sustentáveis, direcionados às necessidades da população local e comprometendo menos o meio ambiente.
Outras respostas recorrentes para o item ’ a ’ foram : a criação de hortas comunitárias, a introdução da permacultura e da agricultura familiar.
Frequentemente, para o item ‘ b ’, os estudantes responderam que seria fundamental promover investimentos na proteção dos recursos hídricos e na saúde pública, também contemplados no padrão de resposta apresentado.
Muitas respostas foram parciais e sem que se observasse preocupação dos estudantes em descreverem / detalharem as duas ações relativas ao item ‘ a ’.
Com apresentação superficial, poucos respondentes associaram as ações sustentáveis apresentadas para a produção de alimentos à redução dos efeitos sobre as mudanças climáticas.
Este foi o maior erro cometido no item ‘ a ’.
O item ‘ b ’ foi o que permitiu o maior número de acertos, em decorrência da amplitude de possibilidades de respostas, decorrentes da abrangência da própria pergunta, e da não necessidade de descrição / detalhamento / associações.
Assim como nas demais questões discursivas, o desempenho dos estudantes foi ruim, neste caso, talvez devido à necessidade de apresentarem quatro ações, sendo maior a complexidade da resposta ao item ‘ a ’.
Algumas deficiências foram detectadas, principalmente na internalização de conhecimentos geográficos mais específicos, a exemplo da temática das mudanças climáticas, tão propagada e debatida pelas mídias.
A dificuldade em fazer determinadas associações, a exemplo da produção de alimentos com a redução dos efeitos \textbf{nocivos} ao clima, mostrou não somente falta de entendimento do que o comando solicitava, mas certa deficiência em explorar conteúdos mais complexos e correlacionáveis, seguramente abordados em algum momento durante a integralização do curso de Licenciatura em Geografia.
Questão 5 – Bacharelado O enunciado da Questão 5, específica do Bacharelado, apresentava um pequeno trecho de um texto sobre planejamento ambiental.
A definição apresentada pelo texto contribuía para que os estudantes respondessem aos itens ‘ a ’ e ‘ b ’ da questão, já que os dois comandos solicitavam explicar o que é planejamento do uso do solo ( rural, no item ‘ a ’, e urbano, no item ‘ b ’ ), e o texto informava que um planejamento ambiental “ obrigatoriamente, deve levar em consideração os aspectos dos sistemas ambientais físicos ( geossistemas ) e dos sistemas socioeconômicos ”.
Já o item ‘ c ’ solicitava citar e explicar um instrumento de planejamento ambiental, o que não estava exemplificado no texto.
O nível de conhecimento exigido pela questão pode ser considerado mediano para os dois primeiros itens e mais alto para o terceiro.
As respostas aos dois primeiros itens, que envolviam explicar o planejamento do uso do solo, poderiam ter estruturas parecidas, e o estudante precisaria mobilizar alguns conceitos.
Para responder ao terceiro item, o estudante precisaria compreender o que é um instrumento de planejamento ambiental, escolher algum dentre os existentes e explicá -lo.
Os conceitos requeridos pela questão faziam parte daqueles listados pelos artigos 6º ( parágrafos XIII e XIV ) e 7º das Diretrizes Curriculares da Área de Geografia, modalidade Bacharelado, no âmbito do Enade, edição 2021.
O padrão de resposta para os itens ‘ a ’ e ‘ b ’ dava as diretrizes gerais que deveriam constar das respostas : que o planejamento do uso do solo é norteado pelo zoneamento, que este define características, potencialidades e restrições do meio físico/biótico e que aspectos afetos ao meio rural para o item ‘ a ’ e urbano para o item ‘ b ’ deveriam ser mencionados.
Para o item ‘ c ’, o padrão de resposta apresentava quatro instrumentos que poderiam ser citados pelos respondentes, com as explicações respectivas, e se referia a outros que poderiam constar das respostas.
As respostas fornecidas pelos estudantes não espelharam o padrão de respostas para os três itens.
No item ‘ c ’, muitos estudantes parafrasearam parte do enunciado da questão, defendendo a importância de se proteger o meio ambiente, sem, necessariamente, responder ao que foi solicitado.
Quando os estudantes respondiam, apenas citavam um instrumento ou uma política de planejamento ambiental, sem explicações.
Um dos tipos de respostas mais frequentes para o item ‘ a ’ foi a afirmação de que planejar o uso do solo rural significa planejar como utilizar o solo de forma a não prejudicar o meio ambiente e, ao mesmo tempo, garantir a qualidade dos produtos agrícolas.
No caso do item ‘ b ’, o planejamento do solo urbano, de acordo com muitos respondentes, representa propor áreas específicas para o desenvolvimento de determinadas atividades urbanas ( zoneamento ), a exemplo de áreas industriais, comerciais, residenciais, etc.
Já nas respostas dadas ao item ‘ c ’, a variedade foi grande, sendo que a maioria dos respondentes citou o instrumento de planejamento ambiental ( muitos desses constantes do padrão de resposta ), mas não explicou o seu significado ou sua importância.
Assim como nas Questões discursivas 3 e 4, o desempenho nessa questão também foi fraco.
As respostas mostraram que os estudantes entendem parcialmente os conceitos relacionados aos dois planejamentos ( rural e urbano ) mas, com relação aos instrumentos de planejamento ambiental, os que responderam, se limitaram a citá -los, parecendo não conseguir explicar o seu funcionamento.
Isso reflete uma falha dos cursos de Geografia em não trabalhar, com mais frequência, temas tão importantes, relacionados ao meio ambiente e à qualidade de vida, que são afetos à Geografia. 6.3.2.7 CONSIDERAÇÕES FINAIS SOBRE AS QUESTÕES DISCURSIVAS DO COMPONENTE DE CONHECIMENTO ESPECÍFICO Foi possível observar, ao longo do processo de correção das Questões Discursivas de Conhecimento Específico de Geografia que os estudantes tiveram um desempenho fraco em todas elas.
Alguns fatores podem ter contribuído para esse quadro : a pandemia da Covid-19 desmotivando a realização da prova ; a forma como algumas questões foram formuladas, com textos motivadores ( enunciados ) relativamente longos ; a natureza de algumas indagações ( de caráter muito abrangente ), o que pode ter levado o estudante, muitas vezes, a não responder efetivamente ao que se pediu, repetindo, como resposta, o enunciado da própria questão.
Somado a isso, tem o fato de alguns estudantes saberem que o que estava sendo efetivamente avaliado eram os cursos e não especificamente o seu desempenho.
Por fim, espera -se que esta e outras avaliações possam servir de base não somente para a reformulação do próprio Enade, como, também, para uma análise do processo de ensino-aprendizagem dos diversos cursos de Geografia ( Licenciatura e Bacharelado ) do país.

% matched lemmas: aprendizado, calça, causalidade, desencadear, nocivo, observação, repor
\end{textsample}
